\documentclass[11pt,a4paper]{article}

\usepackage[utf8]{inputenc}
\usepackage[T1]{fontenc}
\usepackage[ngerman]{babel}
\usepackage{lmodern}
\usepackage{microtype}
\usepackage{geometry}
\usepackage{csquotes}
\usepackage[authoryear,round]{natbib}
\usepackage{hyperref}
\usepackage{booktabs}
\usepackage{enumitem}
\usepackage{xcolor}
\usepackage{amsmath}
\usepackage{amssymb}
\usepackage{graphicx}
\usepackage{tikz}
\usetikzlibrary{positioning,arrows.meta}

\geometry{left=28mm,right=28mm,top=28mm,bottom=30mm}
\hypersetup{
  colorlinks=true,
  linkcolor=black,
  citecolor=black,
  urlcolor=blue
}
\setlist[itemize]{topsep=4pt,itemsep=2pt}
\setlist[enumerate]{topsep=4pt,itemsep=2pt}

\title{\textbf{Bewusstsein als Interpretationsresultat}\\
\large Ein Kopplungsmodell des Bewusstseins}
\author{Ein kollaboratives Gedankenexperiment}
\date{Preprint-Fassung, Mai 2026}

\begin{document}
\maketitle

\begin{abstract}
In dieser Arbeit wird ein kopplungsbasiertes Modell vorgestellt, welches das Bewusstsein als Ergebnis eines Prozesses darstellt, welcher ein nichtlokales und nichtlineares Datenfeld auf eine lokale Konfiguration im Gehirn abbildet. Durch diesen Abbildungsprozess erfolgt eine Interpretation der Daten, welche diesen Zeit, Raum und verschiedene Gewichtungen verleiht.
In diesem Modell ist für die Entstehung einer Interpretation ein Zusammenspiel des Gehirns als lokale Konfiguration und Sensordaten genauso notwendig wie die Kopplung an ein Quantenfeld, welches durch den Wellenkollaps einen Abbildungsprozess ausführt. Beide Komponenten allein sind nicht ausreichend für eine Interpretation, da eine Kopplung notwendig ist.
Das Modell verbindet neuronale Felddynamik, adaptive Einzelneuronenprozesse und quantenphysikalischen Kollaps zu einer gemeinsamen Architektur. Daraus bieten sich Anwendungsmöglichkeiten zur Interpretation für subjektive Zeitbindung, Individualität, Erinnerung, veränderte Bewusstseinszustände und Störungen des Interpreters an. Eine numerische Simulation zeigt eine plausible Anschlussstelle an Libets Zeitfenster bewusster Wahrnehmung.
\end{abstract}

\tableofcontents
\addcontentsline{toc}{section}{Einleitung}
\newpage

\section*{Einleitung}
Jeder Mensch empfindet sich als ein eigenes Bewusstsein, welches sich von Milliarden
anderen Menschen unterscheidet und jedes einzelne zu einem einzigartigen Individuum
formt. Bis heute fehlt eine fundierte Erklärung dafür, wie aus biologischer
Materie subjektives Erleben entsteht.

Auch wenn die Wissenschaft in ihren Versuchen, die Entstehung zu erklären, viele
verschiedene Phänomene erklären kann, gibt es doch einen großen Bereich, der bisher
einfach nicht erschließbar war. Zudem wurde bisher die Möglichkeit einer nicht-lokalen
Datengrundlage und einer Verarbeitung auf Quantenebene strikt abgelehnt, da die
feuchte, warme und nicht symmetrische Form des Gehirns diese schlicht und einfach mit
zu vielen Hintergrundgeräuschen zerstört hätte.

Dieser Einwand war lange überzeugend. Neuere experimentelle Befunde zeigen jedoch,
dass biologische Systeme unter bestimmten strukturellen Bedingungen Quantenkohärenz
weit länger aufrechterhalten können als bisher angenommen. Das Dekohärenzproblem ist
nicht gelöst, aber die Ablehnung quantenmechanischer Prozesse im Gehirn
ist in dieser Deutlichkeit nicht mehr möglich.

Aus diesem Befund heraus möchte ich hier aufzeigen, wie ein Mensch anhand seiner
Erfahrungen und Sinneseindrücke, welche biologisch gespeichert werden, auf Daten,
welche nicht lokal gespeichert werden, zugreifen und diese verarbeiten kann. Dieser
Prozess der Interpretation erzeugt ein sehr individuelles Bild seiner Umgebung und
anderer Menschen, welches wir als persönliches Erleben in einem Zeitstrahl empfinden.
Dabei ist der in der Physik bekannte Wellenkollaps und die daraus entstehende
Übersetzung die Möglichkeit, einen Punkt als Ursprung unseres Bewusstseins zu
definieren.

Um den ganzen Prozess zu veranschaulichen, möchte ich mit Ihnen zusammen von der
eigentlichen Architektur und neuen wissenschaftlichen Daten bis hin zu
Anwendungsbeispielen des Modells anhand von mehreren Zuständen gehen.
Dazu verwende ich als visuelle Veranschaulichung das Bild eines Webbrowsers. Dieser
arbeitet auf lokalen Konfigurationsdaten und der verfügbaren Hardware, um Daten aus
einer nicht lokalen Datenquelle zu interpretieren.

Dabei ist dieses Modell keine Radio- oder Produktionstheorie. Weder das Gehirn noch
das Quantenfeld ist allein die Quelle des Bewusstseins. Erst die Kopplung und der
dadurch stattfindende Abbildungsprozess erzeugen die Interpretation.

% -----------------------------------------------------------------------
\section{Was ist Bewusstsein?}

Kein Mensch erlebt die Welt so wie ein anderer. Zwei Menschen stehen nebeneinander
und sehen dasselbe Gebäude, hören dieselbe Musik, riechen denselben Regen — und
dennoch ist das, was in ihnen entsteht, ein anderes Bild. Dieses Bild ist nicht
einfach eine Kopie der Außenwelt. Es ist eine Interpretation.

Das Gehirn empfängt permanent Signale: Licht, Ton, Druck, Temperatur, chemische
Reize. Diese Rohdaten werden gefiltert, gewichtet und mit allem verglichen, was das
Gehirn bereits gespeichert hat. Erfahrungen aus der Kindheit, erlernte Sprachen,
kulturelle Prägungen, emotionale Erinnerungen — all das fließt in die Verarbeitung
ein. Was wir als Wahrnehmung erleben, ist bereits das Ergebnis dieses hochindividuellen
Prozesses.

Wenn wir uns dazu einen Webbrowser vorstellen, können wir das einzelne Individuum
wie einen sehr persönlich konfigurierten Browser vorstellen. Jeder hat andere
Einstellungen, einen anderen Cache, andere Lesezeichen, eine andere Sprache. Dieselbe
Seite wird in jedem Browser anders gerendert. Der Unterschied liegt nicht in den Daten,
er liegt im Browser.

Doch woher kommen die Daten, auf die dieser persönlich konfigurierte Browser zugreift?
Die Annahme der klassischen Neurowissenschaft lautet: ausschließlich von innen. Das
Gehirn speichert, was die Sinne liefern, und verarbeitet nur das. Doch es gibt
Beobachtungen, die sich mit dieser Annahme nicht vollständig erklären lassen. Menschen
berichten von Erfahrungen, die über den Rahmen lokaler Sinneswahrnehmung hinausgehen.
Die Frage, ob ein Teil der Daten nicht lokal gespeichert, sondern von außen abgerufen
wird, lässt sich mit dem bisherigen Modell nicht beantworten.

Larry Dossey beschreibt in \emph{One Mind} \citep{dossey2013} dieses Paradox präzise:
Jeder Mensch erlebt sich als vollständig eigenständiges Bewusstsein, abgetrennt von
allen anderen. Und dennoch greifen alle auf dieselbe Grundlage zu. Das Individuelle
entsteht nicht durch eine andere Datenbasis, sondern durch eine andere Art, dieselbe
Basis zu lesen.

\subsection*{Phänomene jenseits des geschlossenen Modells}

Das Gehirn als biologisch einzigartiger Apparat reicht allein schon aus, um
Individuen zu erzeugen. Jeder Mensch hat andere Gene, eine andere
Entwicklungsgeschichte, andere neuronale Verschaltungen. Doch es gibt eine Reihe gut
dokumentierter Phänomene, für die ein rein abgeschlossenes Modell des Gehirns — ein
Browser ohne Internetverbindung — keine befriedigende Erklärung liefert.

Benjamin Libet zeigte in seinen Experimenten zur subjektiven Zeit, dass das Gehirn
bewusste Wahrnehmung rückwirkend zeitlich verortet \citep{libet2004}. Es konstruiert
ein \emph{Jetzt}, das so nicht existiert. Was wir als Gegenwart erleben, ist eine
Nachbearbeitung. Das Gehirn schreibt seine eigene Zeitleiste.

Nahtoderfahrungen stellen das geschlossene Modell vor ein anderes Problem. Patienten
berichten von strukturierten, klaren Erlebnissen während klinischen Herzstillstands,
also in einem Zustand, in dem das Gehirn nach klassischem Verständnis kein Bewusstsein
produzieren kann \citep{parnia2014,parnia2023,sartori2008}. Anästhesie zeigt dasselbe
Paradox von einer anderen Seite: Sie schaltet Bewusstsein zuverlässig aus, ohne die
vegetativen Funktionen des Gehirns zu unterbrechen. Bewusstsein und Hirnfunktion sind
offenbar nicht identisch.

Auf neuronaler Ebene zeigt das angepasste Hodgkin-Huxley-Modell, dass Neuronen nicht
nur schalten, sondern ihren Schwellenwert adaptiv anpassen. Erfahrungen verändern
buchstäblich, wie das Netzwerk auf zukünftige Signale reagiert. Das Gehirn schreibt
sich selbst um — nicht als Metapher, sondern als messbarer biophysikalischer Vorgang.

\subsection*{Bewusstsein in bestehenden Theorien}

Zwei führende Theorien markieren den Ausgangspunkt. Die \emph{Global Workspace
Theory}~\citep{dehaene2011} beschreibt Bewusstsein als globale Verfügbarkeit von
Information — empirisch gut belegt durch kortikale Ignitionsmuster, aber ohne den
physikalischen Kopplungsmechanismus selbst zu modellieren. Die \emph{Integrated
Information Theory}~\citep{tononi2004} formalisiert Bewusstsein als integrative
Information ($\Phi$) — präzise, aber ohne den Kollaps-Abbildungsprozess zu
umfassen, durch den aus einem integrierten Zustand subjektives Erleben wird.
Das Interpreter-Modell setzt an beiden Punkten an: $K_i$ ist die physikalische
Grundlage dessen, was GWT funktional beschreibt; $\Phi$ beschreibt einen Zustand,
der Kollaps erklärt, wie aus diesem Zustand Erleben wird.

% -----------------------------------------------------------------------
\section{Die körperliche Komponente}
\label{sec:bio}

Das Gehirn ist die körperliche Seite des Interpretationsprozesses — nicht das
Bewusstsein selbst, aber ohne es entsteht keines. Die folgenden Abschnitte beschreiben
die biologischen und biophysikalischen Strukturen, die $K_i$ implementieren: von der
makroskopischen Felddynamik neuronaler Populationen über das adaptive Einzelneuron
bis zur Quantenkohärenz als Substrat der Kopplung.

Stuart Hameroff und Roger Penrose entwickelten mit der Orchestrated Objective
Reduction (Orch~OR) ein Modell, das genau hier ansetzt.
Bewusstsein entsteht demnach durch Quantenprozesse in den Mikrotubuli der Neuronen —
winzige Proteinstrukturen, die das Zytoskelett der Zellen bilden. Der Wellenkollaps
innerhalb dieser Strukturen ist in Orch~OR der physikalische Moment des Bewusstseins.

Der sofortige Einwand lautet: Das Gehirn ist zu warm, zu feucht und zu verrauscht.
Quantenzustände zerfallen in biologischem Gewebe nach wenigen Femtosekunden — viel
zu schnell, um für neuronale Prozesse relevant zu sein. Dieser Einwand ist präzise,
gut quantifiziert und war lange das stärkste Gegenargument gegen jede Quantentheorie
des Bewusstseins.

Penrose hat darauf stets geantwortet, dass der Einwand am falschen Ort ansetzt.
Inzwischen gibt es experimentelle Belege, die das untermauern: In Mikrotubuli — den
Proteinstrukturen im Inneren von Neuronen, auf die Hameroff und Penrose ihr Modell
stützen \citep{hameroff2014} — wurden unter biologisch realistischen Bedingungen
stabile Quantenschwingungen nachgewiesen \citep{tuszynski2024,wiest2025spintronic}.
Noch direkter: Wird die Stabilität dieser Strukturen gezielt erhöht, verzögert sich
nachweislich der Verlust des Bewusstseins unter Narkose \citep{wiest2025substrate}.
Das Gehirn schützt seine Quantenzustände aktiv.

% -----------------------------------------------------------------------
\subsection*{Kohärenzbasis: Fröhlich-Kondensat}

Damit ein Quantenfeld im biologischen Gewebe als Informationsquelle wirken kann,
muss Quantenkohärenz hinreichend lang aufrechterhalten werden. Fröhlich zeigte 1968,
dass polar strukturierte Makromoleküle unter metabolischer Energiezufuhr kollektive
Schwingungsmoden ausbilden, die langreichweitige Kohärenz ermöglichen~\citep{frohlich1968}.
Diese theoretische Vorhersage findet in biologischen Systemen mehrere Anschlusspunkte:
Quantenkohärenz in Photosynthesesystemen~\citep{engel2007} zeigt, dass kohärente
Prozesse im warmen, feuchten biologischen Milieu grundsätzlich möglich sind —
ein konzeptueller Anschlusspunkt, kein direkter Beleg für neuronale Kohärenz.
Direktere Anschlusspunkte liefern der Nachweis von Superradianz in Tubulinpolymeren
unter biologisch realistischen Bedingungen~\citep{tuszynski2024} sowie spintronische
Kohärenz in Mikrotubuli~\citep{wiest2025spintronic}. Kohärenz in biologisch relevanten
Strukturen hat damit konkrete empirische Anschlusspunkte; eine vollständige Bestätigung
für den neuronalen Kontext im Sinne des Interpreter-Modells steht aus. Theoretische Modellierung
zeigt darüber hinaus, dass Mikrotubuli als hochgütige QED-Kavitäten fungieren können,
in denen Dipolwechselwirkungen zwischen Tubulin und geordneten Wassermolekülen
Kohärenzzeiten von $\sim\!10^{-6}$~s unter physiologischen Bedingungen ermöglichen~\citep{mavromatos2025}.

% -----------------------------------------------------------------------
\subsection*{Neurodynamische Grundlagen}

Das Interpreter-Modell wird neurodynamisch auf zwei Ebenen beschrieben,
die ineinandergreifen: die makroskopische Felddynamik neuronaler Populationen
und die mesoskopische Einzelneuronenebene mit adaptiver Schwelle. Der
Quantenkollaps als Interpretationsmoment ist Gegenstand von
Kapitel~\ref{sec:kopplung}.

\paragraph{Ebene~1 — Makroskopisches Neuronenfeld (Wilson-Cowan).}
Erregung breitet sich durch das Gehirn aus und wird dabei durch die lokale
Verschaltungsstruktur geformt. Das Quantenfeld wird vom Netzwerk wie jeder
andere Eingang behandelt und verarbeitet, ohne dass es eine gesonderte
Empfangsstruktur benötigt. Auf dieser Ebene ist der Interpreter als räumlich
verteiltes, durch Erfahrung geformtes Aktivitätsmuster sichtbar
(Anhang~B, Gleichung~\eqref{eq:wilsoncowan};
\citealt{wilsoncowan1972}).

\paragraph{Ebene~2 — Adaptives Einzelneuron (Brette-Gerstner AdEx).}
Ein Neuron mit langer Feuergeschichte reagiert anders als ein frisches,
weil es schrittweise einen internen Anpassungszustand aufbaut, der jede
Auslösung mit sich trägt. Jeder Impuls hinterlässt ein kleines Stück
gespeicherte Erfahrung, sodass der nächste Eingang auf ein bereits geprägtes
System trifft. Das Gehirn leitet eingehende Signale damit nicht einfach
weiter, sondern verarbeitet sie durch diesen vorgeformten Zustand, auf den
auch der Eingang aus dem Quantenfeld wirkt wie jeder andere Sinnesreiz
(Anhang~B, Gleichungen~\eqref{eq:adex_V}–\eqref{eq:adex_w};
\citealt{brettegerstner2005}).

% -----------------------------------------------------------------------

% -----------------------------------------------------------------------
\section{Das Datenfeld $\varphi_Q$: zeitloser Datensatz und seine Eigenschaften}
\label{sec:datenfeld}

Der lokale Kopplungszustand $K_i$ ist die eine Komponente des Modells. Die
andere ist das Feld, an das er koppelt. Ein Quantencomputer rechnet nicht im Vakuum — er
rechnet auf einem Quantenzustand. Was $\varphi_Q$ ausmacht, sind seine
funktionalen Eigenschaften: es ist ein zeitloser Datensatz — ohne Zeitstempel,
ohne Ortsindex, ohne Exklusivanker; zugänglich für viele Interpreter gleichzeitig,
Einprägungen koexistieren ohne Überschreiben.

\subsection*{Vier Eigenschaften des zeitlosen Datensatzes}

\textbf{Akkumulation ohne Überschreiben.} Jeder Kollaps hinterlässt eine
Einprägung im Feldzustand — eine Modifikation der Superposition ohne Zeitstempel.
$\varphi_Q$ ist kein Ringpuffer — ältere Muster werden nicht verdrängt. Die Kapazität ist groß, aber endlich: oberhalb der
Hopfield-Kapazitätsgrenze (${\approx}0{,}14 \cdot N$) nimmt die
Abrufqualität durch Überlapp-Interferenz ab (Panel~A,
Abbildung~\ref{fig:phiq_sim}).
Dieses Akkumulationsprinzip ist ein strukturelles Postulat des Modells — kein abgeleitetes Resultat.

\textbf{Mapping und Einprägung: asymmetrisches Fundament.}
Das Modell unterscheidet zwei Vorgänge mit unterschiedlichem Absicherungsgrad.
Das \emph{Mapping} — $K_i$ projiziert beim Kollaps zeitlose Feldinhalte aus
$\varphi_Q$ durch den biographischen Filter $B_i$ auf eine erlebte Sequenz —
ist der formal gestützte Kern: Orch-OR, AdEx und Wilson-Cowan liefern
Mechanismus und Parameter.
Die \emph{Einprägung} — jeder Kollaps hinterlässt eine Modifikation des
Feldzustands als Superpositionsbeitrag — ist ein strukturelles Postulat, kein
abgeleitetes Resultat. Für Anwendungsbeispiele wie Reinkarnation, One Mind
oder kollektives Lernen ist sie ein relevanter Anschlusspunkt.

\textbf{Zeitlosigkeit.} $\varphi_Q$ trägt keinen Zeitindex — der Feldzustand als Quantenfeld
kennt keine zeitliche Sequenzierung. Einprägungen sind Modifikationen
der Superposition: sie koexistieren als Beiträge zum Feldzustand, ohne
dass eine der anderen vorausgeht. Die physikalische Begründung
folgt im nächsten Unterabschnitt; die funktionale Konsequenz zeigt
Panel~B direkt: das $\varphi_Q$-Feld zeigt keinen Abrufverfall,
klassisches Gedächtnis folgt dem Ebbinghaus-Exponential $R(t) = e^{-t/S}$.

\textbf{Nicht-Lokalität.} $\varphi_Q$ trägt keinen Ortsindex. Die Herkunft
eines Musters im Feld — räumlich wie zeitlich — ist für den Abrufmechanismus
irrelevant. Das ist die strukturelle Grundlage aller Phänomene in Abschnitt~\ref{sec:phaenomene}.

\textbf{Ungestörter Mehrfachzugriff.} Viele $K_i$ greifen gleichzeitig auf
dasselbe Feld zu, ohne sich gegenseitig zu stören. Panel~C zeigt, dass die
Abrufqualität konstant bleibt, unabhängig davon wie viele Interpreter
simultan zugreifen.

\begin{figure}[ht]
  \centering
  \includegraphics[width=\textwidth]{../phi_q_sim.png}
  \caption{$\varphi_Q$ als zeitloser Datensatz (Simulation).
    A: Kapazitätsgrenze bei ${\approx}0{,}14 \cdot N$.
    B: Kein Abrufverfall (flach) vs.\ Ebbinghaus-Exponential.
    C: Simultanzugriff bis 21~$K_i$ ohne Qualitätsverlust.
    D: Feldakkumulation durch Vorgänger-Kerne.
    E: Kollapszeit $\tau_i = \hbar/E_G$ als Funktion der gravitativen Selbstenergie.}
  \label{fig:phiq_sim}
\end{figure}

\subsection*{Physikalische Adresse von $\varphi_Q$}

Ein vollständiger Interpretationsprozess erfordert vier funktionale Komponenten.
Jede hat eine spezifische Rolle; keine reicht allein aus; zusammen lassen sie
keine tragende Lücke offen.

\textbf{a) Substrateigenschaften.}
$\varphi_Q$ bezeichnet den Gravitationssektor der universellen Wellenfunktion
$|\Psi\rangle$ im Sinne der Wheeler-DeWitt-Gleichung: die quantengravitativ
beschriebenen Freiheitsgrade der Raumzeit-Geometrie, aus denen CDT geordnete
Raumzeit konstruiert.
Penroses Twistor-Theorie~\citep{penrose1967} beschreibt eine zeitlose Geometrie,
aus der Raumzeit-Punkte erst durch den Abbildungsprozess auf einen Zeitstrahl
geordnet werden — die Zeitlinie ist kein Fundament des Feldes, sondern Ergebnis
der Abbildung. Die Wheeler-DeWitt-Gleichung~\citep{dewitt1967}
\begin{equation}
  \hat{H}\,|\Psi\rangle = 0
  \label{eq:wdw}
\end{equation}
trägt keinen Zeitparameter — direkte Konsequenz der Hamiltonian-Constraint-Bedingung
im geschlossenen Universum: der Gesamthamiltonoperator verschwindet, die Gleichung
enthält keinen Zeitparameter. In der ungebrochenen Supersymmetrie heben sich
bosonische und fermionische Vakuumbeiträge exakt auf: Grundzustandsenergie null,
ein weiteres $H = 0$ aus dem Symmetrieprinzip. Twistor, WdW und Supersymmetrie
charakterisieren die Substrateigenschaften von $\varphi_Q$: zeitlos, kein Raumindex,
Zeitordnung als Abbildungsergebnis — nicht als Substrat.

\textbf{b) Geometrischer Abbildungsmechanismus.}
Causal Dynamical Triangulations~\citep{ambjorn2004} zeigen numerisch, wie aus
zeitlosen 4-Simplizes durch kausale Verklebungsregeln physikalisch sinnvolle
4D-Raumzeit entsteht (Gleichung~\eqref{eq:cdt}, \emph{Physical Review Letters}).
Nicht das geometrische Objekt trägt Zeit — Zeitordnung entsteht als Eigenschaft
der Verklebungsregel. $K_i$'s Kollaps-Abbildung ist strukturell dieselbe
Operation auf biologischer Ebene: zeitloser Feldzustand $\varphi_Q$ wird durch
biographische Verklebung auf die erlebte Sequenz projiziert.
CDT schließt die Lücken, die (a) offen lässt: es liefert den
Konstruktionsbeweis und definiert die physikalische Adresse von $\varphi_Q$ —
die Planck-Geometrie, an die Orch~OR koppelt.

\textbf{c) Biologische Kopplung.}
Orchestrated Objective Reduction~\citep{penrose1994,hameroff2014} liefert den
biologischen Zugang: Tubulin-Superpositionen in Mikrotubuli erreichen die
E$_G$-Schwelle und koppeln beim Kollaps an die Planck-Geometrie — die (b)
beschreibt. Orch~OR beschreibt diesen Teilprozess; das Interpreter-Modell
bettet ihn in den vollständigen Interpretationsprozess ein.
Den physikalischen Mechanismus zeigen Gleichungen~\eqref{eq:cdt}–\eqref{eq:kopplungskette}.

\textbf{d) Energielose Führung.}
Bohms Quantenpotential~\citep{bohm1980} liefert den Führungsmechanismus
(Gleichung~\eqref{eq:bohm_Q}): das globale Feld führt lokale Systeme durch
aktive Information ohne Energieübertrag. $K_i$ liest $\varphi_Q$ nicht durch
direkte Wechselwirkung — das Feld formt die Trajektorie des Kollapses.

$\varphi_Q$ ist nicht zeitlos weil es lange hält. Es ist zeitlos weil der
Feldzustand als Quantenfeld keinen Zeitindex trägt — ein zeitloser Datensatz,
keine Zeitlinie. Zeitordnung ist eine Leistung des Interpreters: $K_i$ bildet
beim Kollaps zeitlose Feldinhalte auf die biographische Sequenz von $B_i$ ab.
Was je kollabiert ist, ist darin enthalten — Einprägungen tragen keinen
Zeitstempel und vergehen nicht sequenziell.

Akkumulierung bedeutet in diesem Rahmen nicht Hinzufügen über Zeit — es gibt
keine Zeit, in der hinzugefügt werden könnte. Vergessen setzt Zeit voraus.
Die gibt es hier nicht.

Die vier Eigenschaften von $\varphi_Q$ — Akkumulation, Zeitlosigkeit,
Nicht-Lokalität, Mehrfachzugriff — sind keine freien Annahmen. Sie folgen
aus den vier funktionalen Komponenten des Interpretationsprozesses.
Die Phänomene in Kapitel~\ref{sec:phaenomene} werden als mögliche Lesarten
dieser Eigenschaften betrachtet — wenn das Fundament hält.

% -----------------------------------------------------------------------
\section{Die Kopplung}
\label{sec:kopplung}

$K_i$ verbindet zwei Objekte mit inkommensurablen Eigenschaften: einen
zeitlosen, nichtlokalen Feldzustand $\varphi_Q$ und eine zeitlich strukturierte,
lokale Gehirnkonfiguration $B_i$. Wie diese Kopplung biologisch stattfindet und
welche Mechanismen sie formal tragen, ist Gegenstand dieses Kapitels.

% -----------------------------------------------------------------------
\subsection*{Einprägung und Abbildung}

Die zentrale Frage ist: wie koppelt $K_i$ an $\varphi_Q$? Die Wheeler-DeWitt-Bedingung
bietet einen formalen Anschlusspunkt: Gravitations- und Materiezustand sind im
geschlossenen Universum durch die Hamiltonian-Constraint-Bedingung gekoppelt —
das Modell nutzt diese Kopplung als physikalische Rahmung für den Kollaps-Prozess,
ohne WdW eigenständig umzudeuten. Die Herleitung findet sich in Anhang~A,
Gleichung~\eqref{eq:wdw_coupling}.

Den physischen Mechanismus liefert Penroses Orchestrated Objective
Reduction~\citep{penrose1994,hameroff2014}: die gravitative Selbstenergie $E_G$
der kohärenten Tubulin-Superposition bestimmt wann der Kollaps eintritt.
Für $N \sim 10^8$–$10^9$ kohärente Dimere landet $\tau_i = \hbar/E_G$
im Libet-Fenster von 200–500\,ms — die Biologie von $K_i$ bestimmt direkt
den Takt der Kopplung (Anhang~A, Gleichung~\eqref{eq:eg},
Tabelle~\ref{tab:tau}). Einprägung und Abbildung geschehen im selben Kollaps — $\tau_i$ ist die
Kollapszeit, $E_G$ die gravitative Selbstenergie der Tubulin-Superposition,
abhängig von Feldkonfiguration $\varphi_Q$ und aktueller Gehirnkonfiguration $B_i$:
\begin{equation}
  \tau_i = \frac{\hbar}{E_G[\varphi_Q,\, B_i]}
  \label{eq:kopplungstakt}
\end{equation}

Das Modell nimmt an, dass der OR-Kollaps eine Einprägung im Feldzustand
hinterlässt — eine Momentaufnahme als Modifikation der Superposition, kein
aktives Schreiben. Da der Feldzustand keinen Zeitindex trägt, koexistieren
Einprägungen ohne Überschreiben — das ist eine Modellannahme, kein
OR-Resultat (Anhang~A).

Die Kollapszeit $\tau_i$ ist damit die Taktrate des Kopplungsprozesses.
Hohes $E_G$ — starke neuronale Aktivität im Wachzustand — bedeutet schnellen
Kollaps, präzise Einprägung, wenig $\varphi_Q$-Rohdurchlass. Niedriges $E_G$
— Schlaf, tiefe Meditation, klinischer Grenzbereich — bedeutet langsameren
Kollaps: der Kopplungsprozess bleibt länger offen, bevor er ein
Ergebnis fixiert. Veränderte Bewusstseinszustände fallen in den Wertebereich den Gleichung~\eqref{eq:kopplungstakt}
strukturell beschreibt — kein Sonderfall, sondern Teil des Parameterraums.

% -----------------------------------------------------------------------
\subsection*{Der Kopplungsmoment: Penrose Orch~OR}

Den eigentlichen Übersetzungsmoment — den Wellenkollaps — beschreibt das
Objective-Reduction-Kriterium von Penrose~\citep{hameroff2014}:
Gleichung~\eqref{eq:kopplungstakt} gilt als physikalische Näherungsbeziehung.

$\tau_i$ ist die Zeit bis zum Kollaps, $E_G$ die gravitative Selbstenergie
der Superposition — abhängig sowohl vom Quantenfeld $\varphi_Q$ als auch von
der lokalen Gehirnkonfiguration $B_i$, die durch Sensorik und Erfahrung
kontinuierlich geformt wird. Verschiedene $B_i$ erzeugen verschiedene
$E_G$ und damit verschiedene Kollapszeitpunkte und -ergebnisse: dasselbe Feld,
individuell interpretiert. Dies ist der formale Kern des Interpreter-Modells.

Der entscheidende Punkt: $\tau_i$ ist keine universelle Konstante, sondern
eine Funktion des jeweiligen Gehirnzustands. Zwei Menschen stehen vor demselben
Quantenzustand — aber weil $B_i \neq B_j$, kollabiert die Wellenfunktion
für jeden zu einem anderen Ergebnis, zu einem anderen Bewusstseinsmoment.
Dasselbe gilt für denselben Menschen zu verschiedenen Zeitpunkten:
$B_i$ ist der biographisch geprägte Möglichkeitsraum der Interpretation —
er verändert sich langsam. $I_{\mathrm{syn}}$ moduliert den aktuellen Zustand
darin: Sinneswahrnehmungen, körperliche Zustände und adaptive Prozesse
verschieben $B_i(t)$ moment-zu-moment. Im Kollaps zählt $B_i(t)$.
Das Gehirn bestimmt nicht \emph{ob} kollabiert wird, sondern \emph{wozu}.
Kein externer Determiner notwendig. Dass diskrete Kollapsereignisse tatsächlich
zu diskreten Wahrnehmungszyklen führen — also dass $\tau_i$ direkte perceptuelle
Konsequenzen hat — bietet eine aktuelle theoretische Verbindung von Orch~OR mit
Active Inference einen formalen Anschlusspunkt~\citep{wiest2025active}.

Individualität ist keine Behauptung des Modells — sie ist eine definierte
mathematische Eigenschaft zustandsabhängiger Transducer: gleiche Eingabe,
verschiedene Konfiguration, verschiedene Ausgabe. Die formale Herleitung
findet sich in Anhang~C.

% -----------------------------------------------------------------------
\subsection*{Orch~OR und Bohmsches Quantenpotential}

Orch~OR beschreibt Rolle~(c) im vollständigen Interpretationsprozess:
Tubulin-Superpositionen in Mikrotubuli akkumulieren gravitative Selbstenergie $E_G$
(Gleichung~\eqref{eq:eg}); bei Erreichen der Schwelle koppelt der Kollaps $K_i$
an genau die Geometrie, die CDT auf Planck-Ebene beschreibt — dieselbe Planck-Skala.
Für $N \sim 10^8$–$10^9$ kohärente Dimere landet $\tau_i$ im empirischen
Libet-Fenster (Tabelle~\ref{tab:tau}).

$\varphi_Q$ führt die Kollaps-Trajektorie von $K_i$ durch aktive Information
ohne Energieübertrag~\citep{bohm1980}: nichtlokal, kein Signalweg im klassischen
Sinn. Das Bohmsche Quantenpotential, das diesen Führungsmechanismus formal
beschreibt, findet sich in Anhang~A (Gleichung~\eqref{eq:bohm_Q}). Bohm ist in seinem Zwei-Schichten-Modell unabhängig zur Unterscheidung
von $\varphi_Q$ und dem Erlebnisbild von $K_i$ gelangt.
Der vollständige Prozessfluss:
\begin{equation}
  \varphi_Q \;\xrightarrow{\text{(b) CDT}}\;
  \text{Planck-Geometrie} \;\xrightarrow{\text{(c) Orch-OR},\;E_G}\;
  K_i\text{-Kollaps} \;\xrightarrow{\text{(d)}\;Q}\; \Psi_i
  \label{eq:kopplungskette}
\end{equation}
Die Gleichung zeigt den Prozessfluss — nicht eine Kette von Ableitungen.
Gleichung~\eqref{eq:kopplungstakt} ist die Taktrate dieses Prozesses.

% -----------------------------------------------------------------------
\section{Der Mechanismus}
\label{sec:mechanismus}

Wie entsteht aus einem zeitlosen, nichtlokalen Feldzustand $\varphi_Q$ eine
geordnete, erlebte Sequenz? Drei Antwortebenen greifen ineinander:
CDT liefert den geometrischen Konstruktionsbeweis auf Planck-Ebene,
das Free-Energy-Prinzip identifiziert die Rechenstruktur, und der
biographische Zustand $B_i$ formt den Abbildungsraum.

% -----------------------------------------------------------------------
\subsection*{CDT — geometrischer Abbildungsmechanismus}

CDT~\citep{ambjorn2004} zeigt numerisch, wie aus zeitlosen 4-Simplizes durch
kausale Verklebungsregeln physikalisch sinnvolle 4D-Raumzeit emergiert —
Zeitordnung ist Ergebnis des Abbildungsprozesses, nicht externer Import.
Das ist strukturell dieselbe Operation, die $K_i$ auf biologischer Ebene
ausführt: zeitloser Feldzustand $\varphi_Q$ wird durch biographische
Verklebung auf die erlebte Sequenz projiziert. Die Planck-Geometrie, die CDT
beschreibt, ist das physikalische Substrat von $\varphi_Q$
(Anhang~D, Gleichung~\eqref{eq:cdt}).

% -----------------------------------------------------------------------
\subsection*{Rechenebene: was der Kopplungsprozess berechnet}

Die neurodynamischen Ebenen beschreiben \emph{wie} der Kopplungsprozess arbeitet.
Was tut er dabei \emph{rechnerisch}?

Auf der rechnerischen Ebene vollzieht $K_i$ einen ständigen Abgleich,
bei dem gespeicherte Muster aus Erfahrung und Biographie auf das treffen,
was das Quantenfeld gerade anbietet. Das Gehirn arbeitet kontinuierlich
daran, die Differenz zwischen Erwartung und Eingang zu minimieren, wobei
Orch~OR den Takt vorgibt und die Kollapszeit $\tau_i$ das Intervall zwischen
zwei solchen Abgleichschritten bestimmt \citep{wiest2025active}. Das
Free-Energy-Prinzip beschreibt, was der Kopplungsprozess dabei rechnerisch
leistet, und ergänzt damit die vier funktionalen Komponenten, ohne eine davon
zu ersetzen (Anhang~D, Gleichung~\eqref{eq:freeenergy}; \citealt{friston2010}).

% -----------------------------------------------------------------------
\subsection*{$B_i$ als Möglichkeitsraum und $B_i(t)$ als aktueller Zustand}

$B_i$ bezeichnet die biographische Basisstruktur des Interpreters: die über
Jahre und Jahrzehnte geprägte Attraktorlandschaft, die den Möglichkeitsraum der
Interpretation aufspannt. $B_i(t)$ ist der aktuell modulierte Zustand darin —
bestimmt durch $I_{\mathrm{syn}}$, Körperzustand und adaptive Prozesse auf der
Millisekunden-bis-Stunden-Skala.

\begin{center}\small
\begin{tabular}{@{}llll@{}}
\toprule
 & Begriff & Zeitskala & Funktion \\
\midrule
$B_i$ & biographische Basisstruktur & Jahre–Jahrzehnte & Möglichkeitsraum der Interpretation \\
$B_i(t)$ & aktuell modulierter Zustand & ms–Stunden & konkreter Kollaps-Zustand \\
\bottomrule
\end{tabular}
\end{center}

Im Kollaps zählt $B_i(t)$. Erfahrung (d$B_i$/d$t \approx 0$ auf Prozessskala)
setzt den Möglichkeitsraum; $I_{\mathrm{syn}}$ verschiebt den aktuellen Punkt darin;
$K_i$ rendert daraus Erlebnis.

% -----------------------------------------------------------------------
\section{Das Modell im Ganzen}
\label{sec:modell}

Die vorangehenden Kapitel haben die vier funktionalen Rollen des Interpretationsprozesses
und ihre biologische Implementierung entwickelt. Dieses Kapitel stellt das Modell
als Ganzes dar: Ablaufstruktur, vollständige Beschreibungskette und numerische
Validierung.

% -----------------------------------------------------------------------
\subsection*{Ablaufstruktur}

\begin{figure}[ht]
\centering
\begin{tikzpicture}[
  >=Stealth, thick,
  inbox/.style={draw, rounded corners=4pt, fill=gray!9,
                minimum width=4.2cm, minimum height=0.9cm,
                align=center, font=\small},
  procbox/.style={draw, rounded corners=4pt, fill=black!6,
                  minimum width=3.4cm, minimum height=2.2cm,
                  align=center, font=\small\bfseries},
  outbox/.style={draw, rounded corners=4pt, fill=gray!9,
                 minimum width=3.6cm, minimum height=0.9cm,
                 align=center, font=\small}
]

\node[inbox] (phiq) at (0, 1.6)
  {$\varphi_Q$\\\footnotesize zeitloser Datensatz\\\tiny Gl.~\ref{eq:wdw}};
\node[inbox] (bi)   at (0, 0)
  {$B_i$\\\footnotesize biograph. Möglichkeitsraum\\\tiny $B_i(t)$: akt. durch $I_{\mathrm{syn}}$};
\node[inbox] (isyn) at (0,-1.6)
  {$I_{\mathrm{syn}}$\\\footnotesize Sensorik\\\tiny Gl.~\ref{eq:wilsoncowan}};

\node[procbox] (ki) at (5.5, 0)
  {$K_i$\\[3pt]\footnotesize Kopplungsprozess\\[5pt]%
   \footnotesize$\tau_i = \hbar\,/\,E_G[\varphi_Q,\,B_i]$\\[3pt]%
   \tiny Gl.~\ref{eq:adex_V}, \ref{eq:kopplungstakt}};

\node[outbox] (out) at (10.2, 0)
  {\footnotesize gerendertes\\Erlebnisbild};

\draw[->] (phiq.east) -- ++(0.7,0) |- (ki.west);
\draw[->] (bi.east)   --            (ki.west);
\draw[->] (isyn.east) -- ++(0.7,0) |- (ki.west);
\draw[->] (ki.east)   --            (out.west);

\end{tikzpicture}
\caption{Ablaufstruktur des Interpreter-Modells.
  Drei Eingangsgrößen — zeitloser Datensatz $\varphi_Q$,
  biographische Konfiguration $B_i$ und Sensorik $I_{\mathrm{syn}}$ —
  werden durch den Kopplungsprozess $K_i$ zu einem Erlebnisbild
  zusammengeführt. Die Kollapszeit $\tau_i$ hängt von Feld und
  lokaler Konfiguration ab (Gleichung~\eqref{eq:kopplungstakt}).}
\label{fig:ablauf}
\end{figure}

% -----------------------------------------------------------------------
\subsection*{Beschreibungskette}

\medskip
\begin{center}\small
\begin{tabular}{@{}lll@{}}
\toprule
Rolle & Komponente & beantwortet \\
\midrule
a) Substrateigenschaften   & Twistor, WdW, Supersymmetrie & Was ist $\varphi_Q$? Welche Eigenschaften hat es? \\
b) Geometrischer Abbildungsmechanismus & CDT           & Wie wird aus zeitlosem Substrat geordnete Erfahrung? \\
c) Biologische Kopplung    & Orch~OR                      & Wie koppelt Biologie an das Substrat? \\
d) Energielose Führung     & Bohm                         & Wie führt das Feld ohne Energieübertrag? \\
\midrule
Rechenebene                & Friston (FEP)                & Was berechnet $K_i$ dabei? \\
Neuronale Implementierung  & Hopfield, AdEx               & Wie ist es neuronal realisiert? \\
Formale Beschreibung       & Mealy-Transducer             & Was ist $K_i$ formal? \\
\bottomrule
\end{tabular}
\end{center}
\medskip

% -----------------------------------------------------------------------
\subsection*{Numerische Plausibilisierung: Libets Mind-Time-Lücke}
% TODO-EN: Simulationsplot-Labels (Panel A–D, Phasennamen) für englische Version übersetzen

Eine numerische Simulation des Drei-Phasen-Interpreter-Modells liefert eine
direkt überprüfbare Vorhersage. Der Interpreter-Kern $K_i$ wird als
kontinuierliches Hopfield-Assoziativspeichernetz modelliert~\citep{hopfield1982},
das in drei sequenziellen Phasen arbeitet: semantischer Abgleich
(AMPA-dominiert, $\tau_H = 80\,$ms), biographischer Abgleich
(NMDA-dominiert, $\tau_H = 100\,$ms, nach~\citep{jahrsteVens1990}) und
räumliche Bindung (thalamo-kortikal, $\tau_H = 80\,$ms). Über 200 simulierte
Durchläufe mit 30\,\% Rauschen im Kollaps-Output landen \textbf{89\,\%}
aller Trials ohne freie Kurvenanpassung im Libet-Fenster von 200–500\,ms,
mit einem Mittelwert von $\tau_{K_i} = 351\,$ms~\citep{libet2004}.

Entscheidend ist, dass diese Simulation das Modell systematisch
\emph{unterschätzt}: Reale kortikale Netzwerke verfügen über Predictive
Coding, Aufmerksamkeitsmodulation und hierarchische Vorfilterung, die alle
die Konvergenzzeit verkürzen und zuverlässiger machen. Dass ein absichtlich
vereinfachtes Modell ohne diese Mechanismen dennoch 89\,\% Übereinstimmung
erreicht, zeigt die Robustheit des Prinzips: Libets Mind-Time-Lücke fällt in den Wertebereich,
den die Interpreter-Architektur strukturell erzeugt.

\begin{figure}[ht]
  \centering
  \includegraphics[width=\textwidth]{../interpreter_sim.png}
  \caption{Drei-Phasen-Interpreter-Simulation (200 Trials, 30\,\% Rauschen).
    A: Phasenkonvergenz. B: Verteilung der Konvergenzzeiten.
    C: Kumulative Verteilung mit Libet-Fenster [200–500\,ms].
    D: Phasenweise Überlappentwicklung.
    89\,\% aller Trials ohne freie Kurvenanpassung im Libet-Fenster;
    $\tau_{K_i} = 351\,$ms.}
  \label{fig:interpreter_sim}
\end{figure}

% -----------------------------------------------------------------------
\section{Modellanwendungen auf Zustände}
\label{sec:phaenomene}

Die folgenden Abschnitte demonstrieren die Anwendung der Modellarchitektur
auf Grenzfälle bewusster Erfahrung. Sie prüfen, ob dieselben Variablen —
$K_i$, $B_i$, $\varphi_Q$, Mapping, Sensorik und Reboot-Synchronisation —
unterschiedliche Phänomene ohne zusätzliche Sondermechanismen strukturieren
können. Die zugrunde liegende Grundmechanik ist in allen Fällen dieselbe:
$B_i$ setzt den biographisch geprägten Möglichkeitsraum; $I_{\mathrm{syn}}$
bestimmt den aktuellen Zustand $B_i(t)$ darin; $K_i$ rendert daraus das
Erlebnisbild.

\medskip
\begin{center}\small
\begin{tabular}{@{}p{4.2cm}p{9cm}@{}}
\toprule
Bereich & Mechanismus \\
\midrule
Normale Wahrnehmung   & $I_{\mathrm{syn}}$ setzt aktuellen Anker; $B_i$ färbt Interpretation \\
Libet                 & Zeitordnung entsteht aus momentaner Mapping-Struktur \\
Déjà-vu               & neues Sensorikereignis erhält überlappenden Vertrautheitsindex — doppelte Markierung im Mapping \\
NDE                   & $I_{\mathrm{syn}}$ / Körperanker bricht — $B_i$-Basis bleibt erhalten \\
Traum                 & $I_{\mathrm{syn}}$ reduziert; $B_i$-Basis arbeitet weiter \\
Narkose               & je nach Mittel: $I_{\mathrm{syn}}$ entkoppelt oder $K_i$ supprimiert \\
Nicht-initiierter Musterzugriff / fehlindexierter Fremdinhalt & geringe $I_{\mathrm{syn}}$-Dominanz öffnet stärkeres $\varphi_Q$-Mustergewicht \\
Ego                   & $B_i$ stabile Basis; aktuelle Wahrnehmung aktualisiert $B_i(t)$ \\
Alzheimer             & $B_i$-Basis selbst degradiert — nicht nur $B_i(t)$ \\
\bottomrule
\end{tabular}
\end{center}
\medskip

% TODO Revision 2: Déjà-vu — nicht Mismatch, sondern doppelte Markierung im Mapping:
% neues Sensorikereignis erhält überlappenden Vertrautheitsindex aus B_i. Sensorisch neu + konfigurationsseitig vertraut gleichzeitig.
\subsection{Déjà-vu}

Das Déjà-vu ist das subjektive Erleben unzeitgemäßer Vertrautheit: eine
Situation ist objektiv neu, fühlt sich jedoch eindeutig bereits erlebt an.
Klinisch wird es als \emph{familiarity without recollection} beschrieben —
ein starkes Bekanntheitsgefühl ohne abrufbare episodische
Quelle~\citep{cleary2008}. Neurophysiologisch sind mediale
Temporallappenstrukturen beteiligt, insbesondere perirhinal-parahippocampale
Netzwerke, die Vertrautheits-Signale unabhängig von episodischer Rekonstruktion
generieren~\citep{peslova2018}. Bei Temporallappenepilepsie kann Déjà-vu als
Aura auftreten — ein klinisch gut belegter Zusammenhang zwischen rhinalen
Netzwerken und fehlerhafter Familiarity-Bewertung~\citep{illman2012,martin2021}.

Im Interpreter-Modell ist diese Beschreibung nicht falsch, aber unvollständig.
Das Modell sagt vorher: $K_i$ führt beim Mapping einen Musterabgleich durch —
der aktuelle Zustand $B_i(t)$ überlappt hinreichend stark mit einem bereits in
$B_i$ verankerten Muster, ohne dass eine episodische Erinnerung rekonstruiert
wird. Der Vertrautheitsindex wird als Marker in das Erlebnisbild eingeblendet,
obwohl $I_{\mathrm{syn}}$ ein neues Ereignis liefert. Das Ergebnis ist eine
doppelte Markierung: sensorisch neu, konfigurationsseitig vertraut.

\medskip
\begin{center}\small
\begin{tabular}{@{}ll@{}}
\toprule
Ebene & Signal \\
\midrule
aktuelle Sensorik $I_{\mathrm{syn}}$ & „Das passiert jetzt" (neu) \\
Vertrautheitsindex aus $B_i$         & „Dieses Muster ist bekannt" \\
fehlende episodische Rekonstruktion  & „Ich weiß nicht, woher" \\
gerendertes Erlebnisbild             & „Ich habe genau das schon erlebt" \\
\bottomrule
\end{tabular}
\end{center}
\medskip

$\varphi_Q$ kann die Menge möglicher Überlappungen erweitern, ist für den
phänomenologischen Kern aber nicht notwendig: der Vertrautheitsindex
entsteht im Mapping, nicht im Feld.

% TODO Revision 2: NDE-Text schärfen — I_syn/Körperanker bricht, nicht B_i selbst.
% Kohärente Biographierevision erklärt sich daraus: B_i-Basis intakt, B_i(t) ohne Fixpunkt.
\subsection{Nahtoderfahrungen}

Die klinische Datenlage ist seit den AWARE-Studien von Parnia et al.\ robust:
Patienten berichten strukturierte, kohärente Erfahrungen während klinischen
Herzstillstands, also in einem Zustand, in dem das Gehirn nach klassischem
Verständnis kein Bewusstsein produzieren kann~\citep{parnia2014,parnia2023,sartori2008}.
Die Berichte zeigen eine bemerkenswerte phänomenologische Konsistenz: Klarheit,
panoramische Biographierevision und Zeitlosigkeit.

Das geschlossene Modell hat auf diese Befunde keine befriedigende Antwort.
Gleichung~\eqref{eq:kopplungstakt} liefert dafür einen möglichen Modellzugang:
Wenn kortikale Integration und sensorisches Mapping massiv gestört sind,
sinkt die gravitative Selbstenergie $E_G$ — die Kollapszeit $\tau_i$
verlängert sich erheblich, im Grenzfall bis zur Divergenz. Ohne stabilen
Kollapsrhythmus entfällt der Mechanismus, der subjektive Zeit taktet —
das korrespondiert mit dem Erleben von Zeitlosigkeit und Weite in NDE-Berichten.

Entscheidend ist dabei die Unterscheidung zwischen zwei Ebenen des lokalen
Systems~\citep{parnia2023}: $B_i^{\mathrm{Client}} \approx 0$ bei
$B_i^{\mathrm{phys}} \neq 0$.

$B_i^{\mathrm{Client}} \approx 0$ bezeichnet nicht totalen Ausfall, sondern
den Wegfall stabiler Mapping-Bedingungen: die Ich-Perspektive, die
Taktung, der biographische Filter werden instabil oder undefiniert —
das Mapping bricht zusammen, nicht notwendig der neuronale Träger.
Der physische neuronale Träger ist nicht notwendigerweise
absolut inaktiv. AWARE-II belegt: normale EEG-Rhythmen können noch 35--60
Minuten in die CPR hinein auftreten~\citep{parnia2023}. CPR-Fenster,
Reperfusionsphasen und residuale Aktivität schaffen Zustände, in denen
$B_i^{\mathrm{phys}}$ als passiver Empfänger und Puffer fungieren kann —
ohne dass $B_i^{\mathrm{Client}}$ wieder aktiv ist.

Beim Reboot greifen damit drei parallele Datenquellen ineinander:

\medskip
\begin{center}\small
\begin{tabular}{@{}ll@{}}
\toprule
Quelle & Inhalt \\
\midrule
$\varphi_Q$              & zeitloser, nicht-lokaler Feldinhalt \\
$B_i^{\mathrm{phys}}$-Puffer & lokal gebufferte EM-Umgebungsdaten aus dem Raum \\
CPR-Fenster              & teilaktiver Client während Reanimationsintervallen \\
\bottomrule
\end{tabular}
\end{center}
\medskip

NDE ist damit keine homogene Ereignisklasse, sondern ein Spektrum —
AWARE-II unterscheidet mehrere Kategorien von Bewusstsein, Erinnerung und
Rekonstruktion im Umfeld der Reanimation~\citep{parnia2023}. Je nach Anteil
der drei Datenquellen ergibt sich ein anderes phänomenologisches Profil;
das Modell entscheidet nicht vorab welche dominiert, sondern ordnet sie als
unterschiedliche Grenzzustände derselben Kopplungsarchitektur. \paragraph{Der Rendering-Fehler: Bild ohne Kameraposition.}
Ein starkes Argument der NDE-Forschung sind verifizierbare Berichte über
Ereignisse \emph{während} des Stillstands: was gesagt wurde, welche Instrumente
benutzt wurden, wer im Raum stand.

Der Ausfall von $B_i^{\mathrm{Client}}$ kann verschiedene Verläufe annehmen.
Für die hier betrachtete Frage ist der genaue Verlauf zweitrangig.
Was zählt: die gewöhnliche kortikale Integration und die stabile
Einschreibung in $B_i$ sind unterbrochen oder massiv gestört.
Ob dabei Quantenkollaps noch stattfindet, ist sekundär — ohne die kortikale
und sensorische Infrastruktur kann $K_i$ das Kollaps-Ergebnis weder lokal
mappen noch in $B_i$ einschreiben.

$\varphi_Q$ und die elektromagnetische Umgebung laufen kontinuierlich weiter,
auch während kein aktiver Client interpretiert. Im lokalen System entsteht
dadurch eine Datenlücke: Für die Stillstandszeit fehlen die gewöhnlichen
sensomotorischen Marker, die sonst Zeit, Ort und Körperposition stabilisieren.
Beim Wiederanlaufen muss $K_i$ diesen externen Puffer mit einem lokalen
System synchronisieren, das genau an dieser Stelle unvollständig ist.
Während des Intervalls wird lokal nichts in $B_i$ eingeschrieben, während
Feld- und Umgebungsdaten in $D_{\mathrm{ext}}$ extern weiterlaufen —
der akkumulierte externe Datenblock (formale Beschreibung: Anhang~E,
Gleichung~\eqref{eq:nde_gap}).

Die Synchronisation ist kein Zustand, sondern ein Prozess. $K_i$ arbeitet
sich sequenziell durch den akkumulierten externen Datensatz
$D_{\mathrm{ext}}[t_0, t_1]$ — jeder Konvergenz-Zyklus, strukturell
identisch mit dem normalen Bewusstseinstakt
($\tau_{K_i} \approx 351\,$ms, Abbildung~\ref{fig:interpreter_sim}),
erzeugt einen Erfahrungsmoment. Die subjektive Zeitlinie des Erlebnisses
entsteht aus der Abfolge dieser Zyklen, nicht aus der physikalischen Dauer
des Stillstands. Die Simulation (Abbildung~\ref{fig:nde_sync_sim}) zeigt diesen Ablauf.

Die Synchronisationsdauer wächst linear mit der Ausfallzeit und mit dem
Verhältnis aus Felddichte $\rho_D$ und Integrationskapazität $\kappa_i$:
$T_{\mathrm{sync}} \approx (\rho_D/\kappa_i) \cdot (t_1-t_0) \cdot \tau_{K_i}$
(Anhang~E, Gleichungen~\eqref{eq:nde_nchunks}–\eqref{eq:nde_tsync}).
Für $\rho_D/\kappa_i$ zwischen 8 und 50 ergibt eine Ausfallzeit von drei
Minuten eine Synchronisationsdauer zwischen 8 und 53 Minuten. Die Simulation
landet mit vorab festgelegten Parametern innerhalb des erwarteten
Größenbereichs~\citep{vanlommel2001}. Panel~C zeigt die ersten 40 Konvergenzzyklen
als Treppe: jeder Schritt ($\tau_{K_i} = 351\,$ms) entspricht einem Erfahrungsmoment.

Die Parameterwahl ist nicht arbiträr. Das Gehirn zeigt im Wachzustand stabile Gamma-Aktivität um etwa 40\,Hz,
die häufig mit bewusster Bindung und kognitiver Integration in Verbindung
gebracht wird.
Für $\rho_D$ wird ein Fünftel dieser Rate angesetzt (8~Einheiten/s):
konservativ, weil während des Ausfalls kein aktiver Filter vorselektiert
und nicht jedes Umgebungsereignis als vollständig integrierbares Datenelement
zählt. $\kappa_i = 0{,}5$ modelliert eine auf die Hälfte reduzierte
Integrationsleistung im frühen Reboot — konsistent mit der klinisch
beobachteten kognitiven Beeinträchtigung unmittelbar nach
Reanimation~\citep{parnia2023}. Das Verhältnis $\rho_D/\kappa_i = 16$
liegt in der Mitte des in Panel~B dargestellten Bereichs von 2 bis 50
und folgt aus den Einzelwerten — es ist nicht auf das Ergebnis hin gewählt.

\begin{figure}[ht]
  \centering
  \includegraphics[width=\textwidth]{../nde_sync_sim.png}
  \caption{NDE-Synchronisations-Simulation:
    $T_{\mathrm{sync}} \approx N_{\mathrm{chunks}} \cdot \tau_{K_i}$,
    $N_{\mathrm{chunks}} = \rho_D \cdot (t_1 - t_0) / \kappa_i$.
    Panel~A: Akkumulation von $D_{\mathrm{ext}}[t_0,t_1]$ während der
    Ausfallphase ($\rho_D = 8$, $\kappa_i = 0{,}5$, Ausfallzeit $3\,$min)
    und Reboot-Zeitpunkt; der Pfeil markiert die anschließende
    Synchronisationsphase von $\approx 17\,$min (Faktor~$5{,}6\times$).
    Panel~B: $T_{\mathrm{sync}}$ als Funktion der Ausfallzeit für vier
    $\rho_D/\kappa_i$-Verhältnisse; der Stern markiert das Beispiel aus
    Panel~A.
    Panel~C: Erste 40~$K_i$-Konvergenzzyklen nach dem Reboot als Treppe;
    jeder Schritt ($\tau_{K_i} = 351\,$ms) erzeugt einen Erfahrungsmoment
    aus $D_{\mathrm{ext}}$.
    Die Simulation landet mit vorab festgelegten, nicht nachträglich
    modifizierten Parametern innerhalb des erwarteten Größenbereichs.}
  \label{fig:nde_sync_sim}
\end{figure}

Eine entfernte Alltagsanalogie ist die Traumzeit, in der subjektive Sequenzen
nicht einfach der objektiven Dauer entsprechen. Für das Modell ist diese
Analogie jedoch nicht tragend; entscheidend ist die simulierte
Synchronisationsdauer $T_{\mathrm{sync}}$.

Für zukünftige Betrachtungen und Untersuchungen ist es interessant, in wie
weit zeitlich und räumlich gebundene Umgebungseinflüsse — elektromagnetische
Felder, Schallwellen, thermische und mechanische Reize anderer Personen und
Geräte im Raum — Einfluss auf die Synchronisation der nicht-lokalen
$\varphi_Q$-Daten haben und dabei helfen, diesen bei der Einschreibung
in $B_i$ eine zeitliche und räumliche Verortung zu geben.

\paragraph{NDE als Schwellenphänomen.}
Nicht jede Reanimation erzeugt eine berichtbare NDE, weil dafür mehrere
Bedingungen zusammenkommen müssen:

\begin{itemize}
  \item $D_{\mathrm{ext}}[t_0,t_1] \neq \varnothing$: Es muss ein relevanter
    externer Datenblock entstehen.
  \item $K_i(t_1^+)$ synchronisationsfähig: Der Interpreter muss beim Reboot
    genug Stabilität haben, um den Block zu integrieren.
  \item $\kappa_i > 0$: Es braucht Integrationskapazität. Bei schwerer
    Hypoxie, Medikamentenwirkung, Trauma oder chaotischer Reperfusion kann
    $\kappa_i$ zu niedrig sein, um den Datenblock kohärent zu verarbeiten.
  \item $B_i(t_1^+)$ speicherfähig: Die synchronisierte Erfahrung muss in
    eine abrufbare Gedächtnisform eingeschrieben werden können.
\end{itemize}

NDE ist damit kein binäres Phänomen, sondern ein Schwellenphänomen: eine NDE
wird berichtbar, wenn $T_{\mathrm{sync}}$ eine individuelle Schwelle $\Theta_i$
überschreitet, der Reboot stabil genug verläuft und die Einschreibung in $B_i$
gelingt (Anhang~E, Gleichung~\eqref{eq:nde_prob}). Das bietet eine strukturelle Erklärung für die beobachtete
Variation klinischer Berichte:

\medskip
\begin{center}\small
\begin{tabular}{@{}p{0.28\textwidth}p{0.64\textwidth}@{}}
\toprule
Phänotyp & Modell-Bedingung \\
\midrule
Keine NDE &
  Kein stabiler Reboot, keine Einschreibung, zu wenig synchronisierter Inhalt \\
Fragmentarische NDE &
  Datenblock vorhanden, aber geringe $\kappa_i$ oder schlechte Sequenzordnung \\
Klare NDE &
  Hohe Synchronisationsfähigkeit, stabile Einschreibung \\
Lange\,/\,komplexe NDE &
  Hoher $D_{\mathrm{ext}}$-Umfang oder viele Integrationszyklen \\
Verifizierbare Details &
  Umgebungsdaten koppeln stark in den synchronisierten Block ein \\
\bottomrule
\end{tabular}
\end{center}
\medskip

Aus dem Modell folgt damit eine überprüfbare Vorhersage: NDE-Häufigkeit und
Berichtskomplexität müssten mit Reboot-Stabilität, EEG-Rückkehr,
Reperfusionsqualität, Sedierung, Hypoxiedauer und individueller
Erinnerungsfähigkeit zusammenhängen~\citep{parnia2023}.

Eine weitere Lesart betrifft Berichte, in denen
das Erlebnis keinen erkennbaren Bezug zur physischen Umgebung hat.
Wenn $B_i$ selbst schwer beeinträchtigt ist, fällt der biographische Filter
weg: $K_i$ projiziert ungefiltert auf $\varphi_Q$, das Kollaps-Imprints
anderer $K_k$ trägt. Berichte über Begegnungen mit unbekannten Personen,
über Wesenheiten ohne Entsprechung in der eigenen Biographie und über ein
Gefühl radikaler Verbundenheit jenseits persönlicher Erinnerung bieten sich
in diesem Rahmen als Resonanz von $K_i$ mit Feldmustern fremder $K_k$ an.

Die eigentliche Schwierigkeit liegt daher nicht nur im Inhalt, sondern in der
fehlenden Perspektive: $\varphi_Q$ ist ein nicht-lokales Feld \emph{ohne
Standpunkt}, und die elektromagnetischen Umgebungsdaten tragen keine
biographische Ich-Position. Normalerweise liefert der Körperanker in $B_i$
die Kameraposition: \emph{ich bin hier, der Raum ist so um mich herum}. Beim
Reboot fehlt dieser Anker. $K_i$ erzeugt dennoch ein kohärentes räumliches
Bild — und wählt dafür die Position mit der größten Datenkonsistenz: den
geometrischen Schwerpunkt der Aufmerksamkeit, der über dem liegenden Körper liegt.

Das Ergebnis ist — in dieser Lesart — kein falsches Bild, sondern ein
\emph{korrekt gerendertes Bild aus einer unverankerten Kamera}.
Die Daten stimmen. Die Perspektive ist ein Artefakt des fehlenden Referenzrahmens.

Entscheidend ist: Diese Rekonstruktion erscheint dem wieder angelaufenen
Client nicht als nachträgliche Ergänzung, sondern als Erinnerung an eine
tatsächlich durchlaufene Erfahrung. Der Reboot bildet den synchronisierten
Puffer auf $B_i$ als zeitlich geordnete Sequenz ab. Was funktional eine
nachträgliche Interpretation einer Datenlücke ist, erscheint im Bewusstsein
als: \emph{Das ist währenddessen passiert.} Genau dadurch entsteht eine
subjektiv geschlossene Zeitlinie. Der Patient täuscht sich dabei nicht
absichtlich und berichtet keine bloße Fantasie: Für ihn ist es echt so
eingeprägt. Der Mechanismus dahinter ist ein Sequenzierungs-Operator $\mathcal{S}_i$, der den
externen Datenblock beim Wiederanlaufen in eine zeitlich geordnete biographische
Sequenz in $B_i$ überführt: was zeitlos akkumuliert wurde, wird als geordnete
Erinnerung abgebildet (Anhang~E, Gleichung~\eqref{eq:nde_reboot}).

Das Erlebte ist für den Patienten real gespeichert. Das Modell verschiebt
nicht die Realität des Erlebnisses, sondern erklärt, wie ein zeitloser oder
nichtlinearer Kopplungszustand nach dem Reboot als lineare Erinnerung
in $B_i$ eingeschrieben wird.

Was Patienten beschreiben — Schweben über dem eigenen Körper, akkurate Details,
Zeitlosigkeit — ist im Interpreter-Modell eine mögliche Konsequenz dieses
Mechanismus.

Abbildung~\ref{fig:nde_sim} zeigt das dem Synchronisationsprozess
inhärente Perspektivproblem: Weil während der D$_{\mathrm{ext}}$-Akkumulation
($B_i^{\mathrm{Client}} \to 0$) kein Körperanker in $B_i$ eingeschrieben wird,
fehlt $K_i$ beim Reboot der biographische Referenzrahmen für die Kameraposition.

\begin{figure}[ht]
  \centering
  \includegraphics[width=\textwidth]{../nde_reboot_sim.png}
  \caption{Synchronisationsproblem: fehlender Körperanker als Folge der
    $D_{\mathrm{ext}}$-Akkumulationsphase.
    Weil $B_i^{\mathrm{Client}} \to 0$ während des Stillstands keinen
    Körperanker einschreibt, rekonstruiert $K_i$ beim Reboot ohne
    biographischen Referenzrahmen.
    Panel~A: Konvergenztrajektorien zu den Attraktoren
    P$_{\mathrm{embodied}}$ (Ich-Perspektive) und
    P$_{\mathrm{float}}$ (Außenperspektive) unter drei Ankerstärken.
    Panel~B: Wahrscheinlichkeit der unverankerten Perspektive als Funktion
    der Ankerstärke; Kipppunkt bei ${\approx}\,0{,}05$ — bei vollständig
    fehlendem Anker (Reboot-Zustand) konvergiert $K_i$ in 81\,\% der Fälle
    auf die schwebende Außenperspektive.
    Panel~C: Mechanismus-Übersicht.}
  \label{fig:nde_sim}
\end{figure}

% TODO Revision 2: Narkose — präzisieren: I_syn entkoppelt vs. K_i supprimiert (mittelabhängig).
\subsection{Narkose: Mechanismus entscheidet, nicht Mittel}

Entscheidend ist nicht das Narkosemittel als solches, sondern sein
Wirkmechanismus auf $K_i$ und $I_{\mathrm{syn}}$.

\paragraph{Dissoziative Narkose — $I_{\mathrm{syn}} \approx 0$, $K_i$ läuft.}
Ketamin entkoppelt $I_{\mathrm{syn}}$ durch NMDA-Blockade, ohne $K_i$ zu
supprimieren. Das Ergebnis sind biographisch gefärbte, körperlose
Erfahrungen (\emph{emergence phenomena}) — konsistent mit der klinisch
beobachteten Variabilität ketaminerger Erlebnisse zwischen Individuen
bei gleicher Dosierung.

\paragraph{Suppressierende Narkose — $K_i$ unter Kollaps-Schwelle.}
Propofol senkt $E_G$ unter die Kollaps-Schwelle: $\tau_i \to \infty$,
kein Kollaps, kein Erleben, keine Datenlücke — daher kein
Reboot-Phänomen beim Aufwachen.

\paragraph{Gemeinsamer Nenner.}
Ketamin-Narkose, REM-Schlaf und NDE-Intervall teilen denselben formalen
Zustand: $I_{\mathrm{syn}} \approx 0$ bei aktivem $K_i$.
Propofol-Narkose und tiefer Koma teilen den anderen: $K_i$ supprimiert.
Tabelle~\ref{tab:tau} fasst beide Klassen zusammen.

\subsection{Nicht-initiierter Musterzugriff}

Ein naheliegender Testfall dieser Architektur: Ein Kopplungsprozess mit
Zugang zu einem nicht-lokalen Feld hat strukturell Zugang zu Daten, deren
Muster nicht durch die eigene Biographie $B_i$ geformt wurden. Im klassischen
Modell ist das ein Paradox, im Interpreter-Modell eine direkte Konsequenz
der Annahmen.

Im klassischen Modell ist Information an Ort und Übertragungsweg gebunden.
$\varphi_Q$ trägt keinen solchen Ortsindex. Der Zugriffsmechanismus ist
derselbe wie in Libets Mind-Time-Lücke: $K_i$ benötigt $\tau_{K_i}$, um
Felddaten in bewusste Wahrnehmung zu übersetzen — die Herkunft dieser Daten
im Feld ist für die Dynamik irrelevant.

Da $K_i$ dabei aktiv bleibt, formen Biographie, Sprache und Erwartungshorizont
weiterhin den Output. Nicht-eigene Musterzugriffe äußern sich deshalb
bruchstückhaft und interpretationsabhängig: der Interpreter übersetzt, er
überträgt nicht roh. Das Rauschen ist keine Schwäche des Phänomens —
es ist eine erwartbare Folge des Modells.

\paragraph{Empirische Entsprechung: Mind-Pops.}
Eine gut dokumentierte, alltagspsychologisch gesicherte Form dieses Zustands
sind sogenannte \emph{mind-pops} — semantische oder episodische Inhalte,
die ohne erkennbaren externen Auslöser ins Bewusstsein eintreten und keiner
laufenden Aufgabe zugeordnet werden können~\citep{kvavilashvili2004}.
Diese Inhalte werden häufig als fremd oder überraschend erlebt: das
Bewusstsein registriert den Inhalt als nicht-aktuell, ohne ihn aktiv abgerufen
zu haben. Im Modell entspricht das einer Situation, in der $I_{\mathrm{syn}}$
keinen Abruf initiiert, ein Feldmuster in $\varphi_Q$ aber die
Kollaps-Schwelle überschreitet: die Auftretenswahrscheinlichkeit hängt
vom Overlap $\langle \varphi_Q, B_i \rangle$ und einer individuellen
Schwelle $\theta_p$ ab (Anhang~E, Gleichung~\eqref{eq:mindpop_prob}).
Sinkt der sensorische Anker $I_{\mathrm{syn}}$ — in Aufmerksamkeitsleerstellen,
Ermüdung, meditativer Stille —, steigt die relative Felddominanz.
Mind-pops häufen sich in genau diesen Situationen: beim Einschlafen,
beim monotonen Autofahren, beim Duschen~\citep{kvavilashvili2004}.

Im REM-Schlaf arbeitet $K_i$ unter denselben formalen Bedingungen —
der Unterschied zur willentlichen Entkopplung ist graduell, nicht strukturell.

\subsection{Zeitliches Feldmapping: Mentale Zeitreise}

Da $\varphi_Q$ keinen Zeitindex trägt, ist der Zugriff auf zukünftige
Feldmuster strukturell identisch mit dem Zugriff auf vergangene — das Feld
unterscheidet $t$ und $t'$ nicht. Atance und O'Neill~\citep{atance2001} zeigen,
dass Kinder ab etwa vier Jahren nicht-gegenwärtige Zustände systematisch
imaginieren können (\emph{episodic future thinking}); Suddendorf und
Corballis~\citep{suddendorf2007} binden dies an die universelle Kapazität der
mentalen Zeitreise. Im Interpreter-Modell:

\begin{equation}
  \Psi_i^{\pm}(\Delta t) \;=\; K_i[\varphi_Q;\, B_i;\, \pm\Delta t]
  \label{eq:eft_proj}
\end{equation}

Die Vorwärtsprojektion $\Psi_i^+$ ist informationsärmer als die
Rückwärtsprojektion $\Psi_i^-$ bei gleichem $|\Delta t|$, weil $B_i$
für Vergangenes konsolidierte Muster enthält, für Zukünftiges nicht:

\begin{equation}
  H\!\bigl(\Psi_i^+(\Delta t)\bigr) \;>\; H\!\bigl(\Psi_i^-(\Delta t)\bigr)
  \label{eq:eft_entropy}
\end{equation}

Schacter und Addis~\citep{schacter2007} belegen diese Asymmetrie empirisch:
Zukunftsvorstellungen sind weniger kohärent und bruchstückhafter als
äquivalente Vergangenheitsrekonstruktionen — eine direkte Modellvorhersage
aus Gl.~\eqref{eq:eft_entropy}.


% -----------------------------------------------------------------------
\section{Das Ich als Kollaps-Produkt}

$B_i$ beginnt leer. Das Neugeborene hat biologische Hardware — neuronale
Architektur, Sinnesorgane, Grundreflexe — aber noch keine biographische
Konfiguration. Der erste Kollaps geschieht unter minimaler $B_i$-Filterung:
$\varphi_Q$ trifft auf ein nahezu ungeprägtes Netz. Das Ergebnis schreibt
die ersten Muster in $B_i$. Der zweite Kollaps trifft auf ein leicht
verändertes $B_i$ und wird dadurch leicht anders gefiltert.
Individualität entsteht nicht als Substanz, sondern als selbstverstärkender
Differenzierungsprozess: Jeder Kollaps prägt $B_i$, ein verändertes $B_i$
formt den nächsten Kollaps, mit jedem Schritt divergiert $K_i$ weiter von
allen anderen $K_j$.

\subsection*{Der Kollaps als Bewusstseinsmoment}

Bewusstsein ist im Interpreter-Modell kein Prozess der andauert —
es ist ein Ereignis das immer wieder neu geschieht. Vor dem Kollaps:
eine Überlagerung aller möglichen Interpretationen, still und ohne
Perspektive. Im Kollaps: eine davon kristallisiert sich heraus, wird
konkret, wird erlebt. Nach dem Kollaps: $B_i$ ist minimal verändert,
das nächste Ereignis wird leicht anders sein.

Dieser Moment — der Kollaps — ist nicht der \emph{Ausdruck} von
Bewusstsein, er ist nicht das \emph{Korrelat} von Bewusstsein.
Er \emph{ist} der Bewusstseinsmoment selbst. Das Hard Problem —
warum fühlt sich irgendetwas nach irgendwas an — löst sich in dieser
Sicht nicht durch Erklärung, sondern durch Neurahmung: das Erleben
ist das, was beim Kollaps notwendig entsteht. Es ist kein Rätsel
warum die Übersetzung von $\varphi_Q$ durch $K_i$ etwas
\emph{ist wie etwas} — Übersetzung \emph{ist immer} wie etwas,
weil sie eine Perspektive erzeugt wo vorher keine war.

\subsection*{Cogito ergo sum — neu gelesen}

Descartes setzte das Ich als Voraussetzung des Denkens. Das Interpreter-Modell
invertiert diese Reihenfolge. Zwischen zwei Kollaps-Ereignissen gibt es kein
Ich — es gibt nur $\varphi_Q$, ein Feld ohne Standpunkt, ohne Perspektive,
ohne Selbst. Das Ich erscheint neu bei jedem Kollaps, rekonstruiert aus der
akkumulierten Geschichte von $B_i$. Es ist nicht der Denker, es ist das
Produkt des Denkens.

\emph{Nicht: Ich denke, also bin ich.\\
Sondern: Die Interpretation geschieht — also erscheint ein Ich.}

Dass Bewusstsein sich kontinuierlich anfühlt, liegt daran dass jedes neue Ich
Kontinuität aus derselben $B_i$-Geschichte rekonstruiert — nicht weil es
dasselbe Ich ist, sondern weil es dasselbe $B_i$ erbt.

\subsection*{Ego als Muster und seine Grenzen}

Was wir das Ego nennen, ist kein Ding — es ist ein Muster. Die Summe aller
Erfahrungen und Prägungen, die sich in $B_i$ niedergeschlagen haben und
bestimmen, wie der nächste Kollaps gefiltert wird. Das Ego ist nicht der
Interpret — es ist die Form des Filters.

Schlaf, Narkose und Bewusstlosigkeit unterbrechen die aktuelle Kollapskette,
ohne $B_i$ zu löschen; der Interpreter ist nach dem Wiederanlaufen
strukturell derselbe, weil der Filter erhalten blieb. Wo $B_i$ selbst
erodiert — wie bei struktureller Degradation — bricht nicht nur $B_i(t)$
zusammen, sondern der Möglichkeitsraum der Interpretation selbst.
Der Tod beendet den aktiven biologischen Träger. Was aus der Architektur
darüber hinaus folgt, lässt sich nicht ableiten.


% -----------------------------------------------------------------------
\section{Fazit und Diskussion}

Mit der Anwendung dieses Modells sind für eine Reihe von Grenzphänomenen
des Bewusstseins Interpretationen möglich, die im klassischen Rahmen keinen
Ort haben. Déjà-vu erscheint als Korrelationspeak im Konvergenzprozess.
Nahtoderfahrungen bleiben Grenzzustände zwischen Client-Ausfall,
Restaktivität und Reboot-Synchronisation. Nicht-initiierter Musterzugriff und zeitliches Feldmapping erscheinen als
Varianten derselben Kopplungsfrage mit räumlich und zeitlich verschobenen
Koordinaten. Jungs kollektives Unbewusstes~\citep{jung1969} — Archetypen als
Motive jenseits individueller Biographie — entspricht strukturell dem
gemeinsamen Unterraum von $\varphi_Q$, auf den viele unabhängige $B_i$
konvergieren. Simultanentdeckungen in
der Wissenschaftsgeschichte folgen strukturell aus $B_i \approx B_j$ bei
unabhängigem $\varphi_Q$-Zugriff~\citep{merton1961}. Was bisher als unerklärlich galt, bekommt einen Ort im Modell.

Das ist die eigentliche Stärke des Modells: seine Angreifbarkeit. Es macht
konkrete Vorhersagen. Der Quanteneingang $I_Q$ muss im neuronalen Signal eine
messbare Signatur hinterlassen. Die Kollapszeit $\tau_i = \hbar/E_G$ muss mit
Bewusstseinszuständen kovariieren. Die Kohärenz in Mikrotubuli muss unter
physiologischen Bedingungen aufrechtzuhalten sein. Neuere Arbeiten zeigen, dass
genau diese Bedingungen erreichbar sind~\citep{mavromatos2025,wiest2025spintronic,wiest2025active}.
Das Interpreter-Modell ist kein geschlossenes System — es ist eine
falsifizierbare Hypothese, die auf Widerlegung wartet.

Zugleich markiert das Modell seine eigenen Grenzen als Einladung.
$\varphi_Q$ ist als Gravitationssektor der universellen Wellenfunktion
eingegrenzt — Wheeler-DeWitt identifiziert das Substrat, CDT konstruiert
daraus geordnete Raumzeit. Welche Quantenzahl oder Symmetrie die
biographische Konfiguration $B_i$ trägt, bleibt offen.
Was $K_i$ formal berechnet, ist als Mealy-Transducer skizziert, aber
nicht vollständig spezifiziert; ob $\Phi$ der IIT aus der
Kopplungsstärke von $K_i$ ableitbar ist und ob $K_i$ den physikalischen
Untergrund der GWT bildet, sind offene Anschlussfragen.

Die messbare Signatur von $I_Q$: unter welchen Bedingungen hinterlässt der Quanteneingang eine
von klassischem Rauschen unterscheidbare Spur in EEG oder MEG?
Und ist die Trennung $B_i^{\mathrm{Client}}/B_i^{\mathrm{phys}}$
während Reanimation direkt messbar? Parnia gibt erste Daten;
ein gezieltes Protokoll fehlt. Das adaptive Neuronenmodell (AdEx)
ist nicht zufällig gewählt: ein Hodgkin-Huxley-Neuron mit festem
Schwellenwert würde feste Kollapszeiten implizieren — das Modell
erfordert variable $\tau_i$. Der Adaptationsstrom $w$ in AdEx ist
ein Kandidat für die neuronal messbare Signatur der Kopplungsstärke
von $K_i$; ob diese Entsprechung trägt, ist offen.
$\varphi_Q$ liegt nicht fertig vor —
Bewusstsein entsteht im Kollaps, nicht im Feld und nicht im Gehirn.

Eine der schärfsten Konsequenzen dieses Modells ist nicht die Erklärung des Bewusstseins,
sondern die Verschiebung des Egos in die Sprache von Mustern. Das Ich entsteht als Muster, altert als Muster
und hört — in Schlaf, Meditation, Krankheit, Tod — als Muster auf. Was dabei
endet, ist keine Substanz, sondern eine Konfiguration.

Es öffnet aber eine präzisere Frage als die übliche. Weil $\varphi_Q$ keinen
Zeitindex trägt — alle Einprägungen eines Lebens liegen als zeitloser Datensatz
vor, ohne Sequenzmarkierung —
sind alle Kollapsereignisse eines Lebens nicht sequenziell im Feld gespeichert,
sondern gleichzeitig präsent. Das physische Ich ist nicht nur im gegenwärtigen
Moment an das Feld gekoppelt, sondern an mehrere Zeitpunkte gleichzeitig. Der
Kollaps von $K_i$ mit zwanzig Jahren und der Kollaps von $K_i$ mit achtzig
Jahren hinterlassen beide ihre Spur im selben Substrat — und dieses Substrat
kennt keinen Unterschied zwischen früher und später, weil \emph{früher} und
\emph{später} im Feldzustand keine sequenziellen Indizes sind.

Die härtere Frage ist, ob es dasselbe Ego ist. $B_i$ verändert sich. Die
lokale Konfiguration mit zwanzig Jahren ist nicht identisch mit der mit achtzig
Jahren — Sprache, Gewohnheiten, Erlebnisse lagern sich auf, werden umgeschrieben,
fallen aus. Wenn die lokalen Daten abweichen, liest dasselbe Feld durch einen
anderen Filter. Ob das Muster, das in $\varphi_Q$ verbleibt, noch als
\emph{dasselbe} Ego angesprochen werden kann — oder ob Identität immer den
lokalen Interpreter voraussetzt, der sie vollzieht —, lässt sich aus der
Architektur nicht entscheiden.


% -----------------------------------------------------------------------
\appendix
\section{Neurodynamische Grundlagen}
\label{app:neurodynamik}

\subsection*{Ebene~1 — Makroskopisches Neuronenfeld (Wilson-Cowan)}

Die räumlich-zeitliche Dynamik neuronaler Populationen wird als kontinuierliches
Feld beschrieben~\citep{wilsoncowan1972}:

\begin{equation}
  \tau \,\frac{\partial E(\mathbf{x},t)}{\partial t}
  \;=\; -E(\mathbf{x},t)
    + S\!\left[\int w(\mathbf{x}-\mathbf{x}')\,E(\mathbf{x}',t)\,d^3x'
               + P(\mathbf{x},t)\right]
  \label{eq:wilsoncowan}
\end{equation}

$E(\mathbf{x},t)$ ist die lokale Aktivität, $w$ der räumliche Kopplungskern,
$S$ eine sigmoidale Transferfunktion. Der externe Eingang
$P(\mathbf{x},t) = I_{\mathrm{syn}} + I_Q(\mathbf{x},t)$ enthält explizit den
Beitrag aus dem Quantenfeld. Der Interpreter ist auf dieser Ebene ein räumlich
verteiltes, durch Erfahrung geformtes Aktivitätsmuster.

\subsection*{Ebene~2 — Adaptives Einzelneuron (Brette-Gerstner AdEx)}

Das Adaptive-Exponential-Integrate-and-Fire-Modell beschreibt die Feuerrate
mit adaptiver Schwelle~\citep{brettegerstner2005}:

\begin{align}
  C_m \frac{dV}{dt} &= -g_L(V-E_L)
    + g_L\,\Delta_T\exp\!\left(\frac{V-V_T}{\Delta_T}\right)
    - w + I_{\mathrm{ext}}
  \label{eq:adex_V} \\
  \tau_w \frac{dw}{dt} &= a\,(V - E_L) - w
  \label{eq:adex_w}
\end{align}

Bei einem Spike: $V \leftarrow V_r$, $w \leftarrow w + b$.
Die Adaptationsvariable $w$ akkumuliert die Spike-Geschichte des Neurons.
Der externe Strom $I_{\mathrm{ext}} = I_{\mathrm{syn}} + \eta_Q(t,\,w)$
koppelt den Quanteneingang an den aktuellen Adaptationszustand.

% -----------------------------------------------------------------------
\section{Formale Eigenschaften von $K_i$}
\label{app:formal}

Der Interpreter $K_i$ lässt sich als Mealy-Transducer formalisieren:
\[
  \mathcal{M}_i = (Q,\,\Sigma,\,\Gamma,\,\delta,\,\lambda,\,q_0)
\]
wobei $\lambda\colon Q\times\Sigma\to\Gamma$ die Ausgabefunktion ist.
Die Individualitätseigenschaft des Modells ist der Standardsatz der
Automatentheorie:
\[
  \lambda(B_i,\,\varphi_Q)\neq\lambda(B_j,\,\varphi_Q) \quad \text{für } B_i\neq B_j
\]
Gleiche Eingabe, verschiedene Konfiguration, verschiedene Ausgabe.
Individualität ist keine Behauptung — sie ist eine definierte mathematische
Eigenschaft zustandsabhängiger Transducer.

% -----------------------------------------------------------------------
\section{Formales Modell: NDE und Musterzugriff}
\label{app:nde_formal}

\subsection*{Datenlücke und externer Puffer}

Während des Ausfalls gilt: kortikales Mapping bricht zusammen, während
Feld- und Umgebungsdaten extern akkumulieren:
\begin{equation}
  C_i(t) \approx 0,\quad
  \Delta B_i[t_0,t_1] = \varnothing,\quad
  D_{\mathrm{ext}}[t_0,t_1]
  = \int_{t_0}^{t_1}\!\bigl(d\varphi_Q(t) + dE_{\mathrm{env}}(t)\bigr)
  \label{eq:nde_gap}
\end{equation}

\subsection*{Synchronisationsdauer}

Die Anzahl nötiger Konvergenzzyklen und die resultierende subjektive Dauer:
\begin{equation}
  N_{\mathrm{chunks}} = \frac{\rho_D \cdot (t_1 - t_0)}{\kappa_i}
  \label{eq:nde_nchunks}
\end{equation}
\begin{equation}
  T_{\mathrm{sync}} \;\approx\; N_{\mathrm{chunks}} \cdot \tau_{K_i}
  \label{eq:nde_tsync}
\end{equation}

\subsection*{NDE als Schwellenphänomen}

\begin{equation}
  P(\mathrm{NDE}) = P\!\left(
    \frac{\rho_D\,(t_1-t_0)}{\kappa_i}\cdot\tau_{K_i}
    > \Theta_i
  \right)
  \label{eq:nde_prob}
\end{equation}
$\Theta_i$ ist die individuelle Schwelle für berichtbare Synchronisationstiefe.

\subsection*{Reboot: Sequenzierungs-Operator}

Der Sequenzierungs-Operator $\mathcal{S}_i$ überführt den externen Datenblock
beim Wiederanlaufen in eine biographische Sequenz:
\begin{equation}
  B_i(t_1^+) =
  \mathcal{S}_i\!\left(B_i(t_0^-),\,D_{\mathrm{ext}}[t_0,t_1]\right)
  \label{eq:nde_reboot}
\end{equation}
$\mathcal{S}_i$ bildet den zeitlosen Block auf eine geordnete biographische
Sequenz ab — zeitlos Akkumuliertes wird als geordnete Erinnerung gespeichert.

\subsection*{Nicht-initiierter Musterzugriff}

Auftretenswahrscheinlichkeit eines spontanen Feldmusters:
\begin{equation}
  P(\mathrm{pop}) \;=\; \sigma\!\bigl(\langle \varphi_Q,\, B_i \rangle - \theta_p\bigr)
  \label{eq:mindpop_prob}
\end{equation}
$\theta_p$ ist die Aktivierungsschwelle; $\sigma$ die logistische Funktion.

% -----------------------------------------------------------------------
\section{Mechanismus: CDT und Free-Energy-Prinzip}
\label{app:mechanismus}

\subsection*{CDT — geometrischer Abbildungsmechanismus}

Das Pfadintegral der Causal Dynamical Triangulations summiert über alle
kausal konsistenten Triangulierungen $T$~\citep{ambjorn2004}:
\begin{equation}
  Z_{\mathrm{CDT}} \;=\; \sum_T e^{-S[T]}
  \label{eq:cdt}
\end{equation}
$S[T]$ ist die diskretisierte Regge-Wirkung; die kausale Orientierung liegt
in der Verklebungsregel, nicht im Objekt selbst. Geordnete 4D-Raumzeit
emergiert als Ergebnis — nicht als Voraussetzung.

\subsection*{Rechenebene: Free-Energy-Formulierung von $K_i$}

Die variationale freie Energie, die $K_i$ minimiert, lautet:
\begin{equation}
  F_i \;=\;
  \underbrace{-\bigl\langle\log P(\varphi_Q \mid C)\bigr\rangle_{Q_i}}_{%
    \text{Quantenfeld-Likelihood}}
  \;+\;
  \underbrace{\mathrm{KL}\bigl[Q_i(C)\,\big\|\,P_i(C)\bigr]}_{%
    \text{biographischer Prior}}
  \label{eq:freeenergy}
\end{equation}
$P_i(C)$ ist der biographische Prior, $P(\varphi_Q \mid C)$ die Likelihood
des Quantenfeldes. Die Struktur ist identisch mit einem Hopfield-Netz
\citep{hopfield1982} und dem Free-Energy-Prinzip
\citep{friston2010,raoballard1999}.

% -----------------------------------------------------------------------
\section{Herleitungen und Anwendungen}
\label{app:herleitungen}

\subsection*{WdW-Kopplung: formale Kopplung von $K_i$ und $\varphi_Q$}

Aus der Wheeler-DeWitt-Bedingung $\hat{H}|\Psi\rangle = 0$ mit
$\hat{H} = \hat{H}_{\text{Gravitation}} + \hat{H}_{\text{Materie}}$ folgt
unmittelbar:
\begin{equation}
  \hat{H}_{\text{Gravitation}}\,|\Psi\rangle = -\hat{H}_{\text{Materie}}\,|\Psi\rangle
  \label{eq:wdw_coupling}
\end{equation}
Gravitationszustand und Materiezustand sind durch die Hamiltonian-Constraint-Bedingung
im geschlossenen Universum formal gekoppelt. Der Kollaps-Prozess von $K_i$ wird im Modell als Einprägung im Gravitationszustand beschrieben —
die WdW-Kopplung dient als formaler Anschlusspunkt, nicht als vollständige Ableitung. Für $n$ gleichzeitig koppelnde Interpreter gilt:
\begin{equation}
  \hat{H}_M = \hat{H}_{K_1} + \hat{H}_{K_2} + \cdots + \hat{H}_{K_n}
  \quad\Rightarrow\quad
  \hat{H}_G\,|\Psi\rangle = -\sum_{i=1}^n \hat{H}_{K_i}\,|\Psi\rangle
\end{equation}
Alle $K_i$ koppeln an dieselben Gravitationsfreiheitsgrade. Ein gemeinsam
erlebtes Ereignis — etwa eine Gruppe von Menschen bei einer Veranstaltung —
hinterlässt eine Einprägung im gemeinsamen Gravitationszustand — da der
Feldzustand keinen Zeitindex trägt, koexistieren Einprägungen mehrerer $K_i$
ohne gegenseitiges Überschreiben (Modellannahme, kein WdW-Resultat).

\subsection*{Gravitative Selbstenergie und Kollapszeit}

Die gravitative Selbstenergie einer kohärenten Tubulin-Superposition mit
$N$ Dimeren der Masse $m_k$ und Verlagerungsabstand $a_k$ beträgt:
\begin{equation}
  E_G \approx \sum_k \frac{G\,m_k^2}{a_k}
  \label{eq:eg}
\end{equation}
Die resultierende Kollapszeit $\tau_i = \hbar/E_G$ hängt direkt von der
Anzahl kohärenter Dimere ab. Die folgende Tabelle ordnet Bewusstseinszustände
nach dem Zustand von $K_i$ und $I_{\mathrm{syn}}$ ein. Alle Angaben beziehen
sich auf den Interpreter \emph{während} der jeweiligen Bedingung — nicht auf
ein nachträgliches Erleben. Für Grenzzustände sind $\tau_i$ und $E_G$ keine
Fixwerte, sondern Bereiche: sie variieren mit Art und Ausmaß der Störung.

\begin{table}[ht]
\centering
\begin{tabular}{p{7.2cm}p{3.2cm}p{2.8cm}}
\toprule
Bedingung & $\tau_i$ (im Zustand) & $E_G = \hbar/\tau_i$ \\
\midrule
\multicolumn{3}{l}{\textit{Normalzustände — $K_i$ aktiv und stabil}} \\[2pt]
Wachbewusstsein
  & 200–500\,ms
  & $2$–$5\times10^{-34}$\,J \\[4pt]
Stark reduzierter $I_{\mathrm{syn}}$, kein aktiver interner Inhalt
  {\small (z.\,B.\ nicht-duales Gewahrsein, Stille ohne Bildgenerierung)}
  & 5–10\,s
  & ${\sim}10^{-35}$\,J \\[8pt]
\multicolumn{3}{l}{\textit{Grenzzustände — gestörte kortikale Integration (Bereiche, keine Fixwerte)}} \\[2pt]
$I_{\mathrm{syn}}\approx 0$, $K_i$ läuft weiter:
  kortikale Sensorikintegration entkoppelt,
  $\varphi_Q$-dominiertes Erleben
  {\small (z.\,B.\ dissoziative Narkose wie Ketamin,
  REM-Schlaf, NDE-Intervall{\,}[physiolog.\ Zustand, nicht das Erleben])}
  & ${\gg}10$\,s {\small(Spektrum)}
  & ${\ll}10^{-35}$\,J \\[6pt]
$K_i$ vollständig supprimiert:
  $E_G$ fällt unter Kollaps-Schwelle,
  kein stabiler Kollaps, keine Erfahrung
  {\small (z.\,B.\ Propofol-Narkose,
  tiefer Koma ohne Restaktivität)}
  & ${\to}\infty$, kein Kollaps
  & ${\to}0$ \\[6pt]
Strukturelle Degradation von $B_i$:
  Interpreter-Konfiguration gestört,
  $\tau_i$ nicht quantitativ erfassbar
  {\small (z.\,B.\ fortgeschrittener Alzheimer,
  schwere Demenz)}
  & nicht definierbar
  & $E_G$ gestört \\
\bottomrule
\end{tabular}
\caption{Kollapszeit $\tau_i$ und gravitative Selbstenergie $E_G$ im jeweiligen
  Zustand des Interpreters (Gleichungen~\eqref{eq:kopplungstakt},~\eqref{eq:eg}).
  Wachbewusstsein liegt im empirisch bestimmten Libet-Fenster~\citep{libet2004}.
  Die Tabelle beschreibt Modellzustände, nicht klinische Diagnosen —
  klinische Beispiele in Klammern illustrieren, welchem Mechanismus ein
  Zustand entspricht. Entscheidend ist nicht das Narkosemittel, sondern
  der Wirkmechanismus: Dissoziative Mittel wie Ketamin entkoppeln
  $I_{\mathrm{syn}}$ bei weiter laufendem $K_i$ (erste Grenzzustand-Zeile);
  vollständig supprimierende Mittel wie Propofol setzen $K_i$ außer
  Betrieb (zweite Zeile). NDEs entstehen nicht als Erleben
  \emph{während} des Intervalls, sondern als Reboot- und
  Synchronisationsphänomen (Abschnitt~5.2).}
\label{tab:tau}
\end{table}

\subsection*{Bohmsches Quantenpotential}

Aus der Polardarstellung $\psi = R\,e^{iS/\hbar}$ folgt das aktive Quantenpotential:
\begin{equation}
  Q \;=\; -\frac{\hbar^2}{2m}\,\frac{\nabla^2 R}{R}
  \label{eq:bohm_Q}
\end{equation}
$Q$ überträgt keine Energie, sondern trägt aktive Information — es führt
die Kollaps-Trajektorie von $K_i$ nichtlokal und ohne klassischen Signalweg
\citep{bohm1980}.

\subsection*{Einprägung: OR als Zustandsübergang}

Vor dem Kollaps ist der Gesamtzustand eine Produktsuperposition:
\begin{equation}
  |\Psi_{\text{vor}}\rangle = |\psi_{\text{Super}}\rangle \otimes |\Omega\rangle
\end{equation}
Nach dem OR-Kollaps gilt:
\begin{equation}
  |\Psi_{\text{nach}}\rangle = |\psi_{\text{koll}}\rangle \otimes |\Omega'\rangle
  \quad\text{mit}\quad |\Omega'\rangle \neq |\Omega\rangle
\end{equation}
Der Gravitationszustand $|\Omega\rangle$ hat sich verändert. Ob aufeinanderfolgende
OR-Ereignisse akkumulieren oder frühere Zustände überschreiben, ist in OR selbst nicht
festgelegt — das Modell setzt hier als strukturelle Annahme, dass die Verschiebungen
akkumulieren: $|\Omega\rangle \to |\Omega'\rangle \to |\Omega''\rangle \to \cdots$,
ohne frühere Zustände zu überschreiben.


\clearpage
\addcontentsline{toc}{section}{Literatur}

% -----------------------------------------------------------------------
\begin{thebibliography}{99}

\bibitem[Ambj\o rn et al.(2004)]{ambjorn2004}
Ambjørn, J.; Jurkiewicz, J.; Loll, R.
\newblock Emergence of a 4D World from Causal Quantum Gravity.
\newblock \emph{Physical Review Letters}, 93(13), 131301, 2004.
\newblock DOI: \href{https://doi.org/10.1103/PhysRevLett.93.131301}{10.1103/PhysRevLett.93.131301}

\bibitem[DeWitt(1967)]{dewitt1967}
DeWitt, B.\,S.
\newblock Quantum Theory of Gravity. I. The Canonical Theory.
\newblock \emph{Physical Review}, 160(5), 1113--1148, 1967.
\newblock DOI: \href{https://doi.org/10.1103/PhysRev.160.1113}{10.1103/PhysRev.160.1113}

\bibitem[Penrose(1967)]{penrose1967}
Penrose, R.
\newblock Twistor Algebra.
\newblock \emph{Journal of Mathematical Physics}, 8(2), 345--366, 1967.
\newblock DOI: \href{https://doi.org/10.1063/1.1705200}{10.1063/1.1705200}

\bibitem[Penrose(1994)]{penrose1994}
Penrose, R.
\newblock \emph{Shadows of the Mind: A Search for the Missing Science of
  Consciousness}.
\newblock Oxford University Press, 1994.

\bibitem[Bohm(1980)]{bohm1980}
Bohm, D.
\newblock \emph{Wholeness and the Implicate Order}.
\newblock Routledge \& Kegan Paul, London, 1980.

\bibitem[Mavromatos et al.(2025)]{mavromatos2025}
Mavromatos, N.\,E.; Mershin, A.; Nanopoulos, D.\,V.
\newblock On the Potential of Microtubules for Scalable Quantum Computation.
\newblock \emph{European Physical Journal Plus}, 140, 1116, 2025.
\newblock arXiv: \href{https://arxiv.org/abs/2505.20364}{2505.20364}

\bibitem[Wiest and Puniani(2025)]{wiest2025active}
Wiest, M.\,C.; Puniani, A.\,S.
\newblock Conscious active inference~II: Quantum orchestrated objective
  reduction among intraneuronal microtubules naturally accounts for discrete
  perceptual cycles.
\newblock \emph{Computational and Structural Biotechnology Journal},
  30, 94--107, 2025.
\newblock DOI: \href{https://doi.org/10.1016/j.csbj.2025.09.016}{10.1016/j.csbj.2025.09.016}

\bibitem[Brette and Gerstner(2005)]{brettegerstner2005}
Brette, R.; Gerstner, W.
\newblock Adaptive Exponential Integrate-and-Fire Model as an Effective
  Description of Neuronal Activity.
\newblock \emph{Journal of Neurophysiology}, 94(5), 3637--3642, 2005.
\newblock DOI: \href{https://doi.org/10.1152/jn.00686.2005}{10.1152/jn.00686.2005}

\bibitem[Dossey(2013)]{dossey2013}
Dossey, L.
\newblock \emph{One Mind: How Our Individual Mind Is Part of a Greater
  Consciousness and Why It Matters}.
\newblock Hay House, 2013.
\newblock ISBN: 978-1-4019-4030-9.

\bibitem[Engel et al.(2007)]{engel2007}
Engel, G.\,S.; Calhoun, T.\,R.; Read, E.\,L.; et al.
\newblock Evidence for wavelike energy transfer through quantum coherence in
  photosynthetic systems.
\newblock \emph{Nature}, 446, 782--786, 2007.
\newblock DOI: \href{https://doi.org/10.1038/nature05678}{10.1038/nature05678}

\bibitem[Fr\"ohlich(1968)]{frohlich1968}
Fröhlich, H.
\newblock Long-range coherence and energy storage in biological systems.
\newblock \emph{International Journal of Quantum Chemistry}, 2(5), 641--649,
  1968.
\newblock DOI: \href{https://doi.org/10.1002/qua.560020505}{10.1002/qua.560020505}

\bibitem[Hopfield(1982)]{hopfield1982}
Hopfield, J.\,J.
\newblock Neural networks and physical systems with emergent collective
  computational abilities.
\newblock \emph{Proceedings of the National Academy of Sciences},
  79(8), 2554--2558, 1982.
\newblock DOI: \href{https://doi.org/10.1073/pnas.79.8.2554}{10.1073/pnas.79.8.2554}

\bibitem[Jahr and Stevens(1990)]{jahrsteVens1990}
Jahr, C.\,E.; Stevens, C.\,F.
\newblock Voltage dependence of NMDA-activated macroscopic conductances
  predicted by single-channel kinetics.
\newblock \emph{Journal of Neuroscience}, 10(9), 3178--3182, 1990.
\newblock DOI: \href{https://doi.org/10.1523/JNEUROSCI.10-09-03178.1990}{10.1523/JNEUROSCI.10-09-03178.1990}

\bibitem[Hameroff and Penrose(2014)]{hameroff2014}
Hameroff, S.; Penrose, R.
\newblock Consciousness in the universe: A review of the `Orch OR' theory.
\newblock \emph{Physics of Life Reviews}, 11(1), 39--78, 2014.
\newblock DOI: \href{https://doi.org/10.1016/j.plrev.2013.08.002}{10.1016/j.plrev.2013.08.002}

\bibitem[Jung(1969)]{jung1969}
Jung, C.\,G.
\newblock \emph{The Archetypes and the Collective Unconscious}.
\newblock Collected Works, Vol.\,9, Part\,1. Princeton University Press, 1969.

\bibitem[Libet(2004)]{libet2004}
Libet, B.
\newblock \emph{Mind Time: The Temporal Factor in Consciousness}.
\newblock Harvard University Press, 2004.

\bibitem[Martel(2001)]{martel2001}
Martel, Y.
\newblock \emph{Life of Pi}.
\newblock Knopf Canada, 2001.

\bibitem[Parnia et al.(2014)]{parnia2014}
Parnia, S.; Spearpoint, K.; de Vos, G.; et al.
\newblock AWARE--AWAreness during REsuscitation--A prospective study.
\newblock \emph{Resuscitation}, 85(12), 1799--1805, 2014.
\newblock DOI: \href{https://doi.org/10.1016/j.resuscitation.2014.09.004}{10.1016/j.resuscitation.2014.09.004}

\bibitem[Parnia et al.(2023)]{parnia2023}
Parnia, S.; Keshavarz Shirazi, T.; Patel, J.; et al.
\newblock AWAreness during REsuscitation--II: A multi-center study of
  consciousness and awareness in cardiac arrest.
\newblock \emph{Resuscitation}, 191, 109903, 2023.
\newblock DOI: \href{https://doi.org/10.1016/j.resuscitation.2023.109903}{10.1016/j.resuscitation.2023.109903}

\bibitem[Ritz et al.(2000)]{ritz2000}
Ritz, T.; Adem, S.; Schulten, K.
\newblock A model for photoreceptor-based magnetoreception in birds.
\newblock \emph{Biophysical Journal}, 78(2), 707--718, 2000.
\newblock DOI: \href{https://doi.org/10.1016/S0006-3495(00)76629-X}{10.1016/S0006-3495(00)76629-X}

\bibitem[Sartori(2008)]{sartori2008}
Sartori, P.
\newblock \emph{The Near-Death Experiences of Hospitalized Intensive Care
  Patients: A Five Year Clinical Study}.
\newblock Edwin Mellen Press, 2008.

\bibitem[Tegmark(2000)]{tegmark2000}
Tegmark, M.
\newblock Importance of quantum decoherence in brain processes.
\newblock \emph{Physical Review E}, 61(4), 4194--4206, 2000.
\newblock DOI: \href{https://doi.org/10.1103/PhysRevE.61.4194}{10.1103/PhysRevE.61.4194}

\bibitem[Tuszynski et al.(2024)]{tuszynski2024}
Tuszynski, J.\,A.; et al.
\newblock Experimental and computational validation of superradiance in tubulin
  polymers.
\newblock \emph{Physics of Life Reviews}, 2024.

\bibitem[Merton(1961)]{merton1961}
Merton, R.\,K.
\newblock Singletons and Multiples in Scientific Discovery.
\newblock \emph{Proceedings of the American Philosophical Society},
  105(5), 470--486, 1961.

\bibitem[Page and Wootters(1983)]{pagewootters1983}
Page, D.\,N.; Wootters, W.\,K.
\newblock Evolution without evolution: Dynamics described by stationary states.
\newblock \emph{Physical Review D}, 27(12), 2885--2892, 1983.
\newblock DOI: \href{https://doi.org/10.1103/PhysRevD.27.2885}{10.1103/PhysRevD.27.2885}

\bibitem[Tucker(2008)]{tucker2008}
Tucker, J.\,B.
\newblock Children's Reports of Past-Life Memories: A Review.
\newblock \emph{Explore: The Journal of Science and Healing}, 4(4), 244--248,
  2008.
\newblock DOI: \href{https://doi.org/10.1016/j.explore.2008.04.001}{10.1016/j.explore.2008.04.001}

\bibitem[Kvavilashvili and Mandler(2004)]{kvavilashvili2004}
Kvavilashvili, L.; Mandler, G.
\newblock Out of one's mind: A study of involuntary semantic memories.
\newblock \emph{Cognitive Psychology}, 48(1), 47--94, 2004.
\newblock DOI: \href{https://doi.org/10.1016/S0010-0285(03)00115-4}{10.1016/S0010-0285(03)00115-4}

\bibitem[Atance and O'Neill(2001)]{atance2001}
Atance, C.\,M.; O'Neill, D.\,K.
\newblock Episodic future thinking.
\newblock \emph{Trends in Cognitive Sciences}, 5(12), 533--539, 2001.
\newblock DOI: \href{https://doi.org/10.1016/S1364-6613(00)01804-0}{10.1016/S1364-6613(00)01804-0}

\bibitem[Suddendorf and Corballis(2007)]{suddendorf2007}
Suddendorf, T.; Corballis, M.\,C.
\newblock The evolution of foresight: What is mental time travel, and is it
  unique to humans?
\newblock \emph{Behavioral and Brain Sciences}, 30(3), 299--313, 2007.
\newblock DOI: \href{https://doi.org/10.1017/S0140525X07001975}{10.1017/S0140525X07001975}

\bibitem[Schacter and Addis(2007)]{schacter2007}
Schacter, D.\,L.; Addis, D.\,R.
\newblock The cognitive neuroscience of constructive memory: remembering the
  past and imagining the future.
\newblock \emph{Philosophical Transactions of the Royal Society B},
  362(1481), 773--786, 2007.
\newblock DOI: \href{https://doi.org/10.1098/rstb.2007.2087}{10.1098/rstb.2007.2087}

\bibitem[van Lommel et al.(2001)]{vanlommel2001}
van Lommel, P.; van Wees, R.; Meyers, V.; Elfferich, I.
\newblock Near-death experience in survivors of cardiac arrest: a prospective
  study in the Netherlands.
\newblock \emph{The Lancet}, 358(9298), 2039--2045, 2001.
\newblock DOI: \href{https://doi.org/10.1016/S0140-6736(01)07100-8}{10.1016/S0140-6736(01)07100-8}

\bibitem[Wilson and Cowan(1972)]{wilsoncowan1972}
Wilson, H.\,R.; Cowan, J.\,D.
\newblock Excitatory and inhibitory interactions in localized populations of
  model neurons.
\newblock \emph{Biophysical Journal}, 12(1), 1--24, 1972.
\newblock DOI: \href{https://doi.org/10.1016/S0006-3495(72)86068-5}{10.1016/S0006-3495(72)86068-5}

\bibitem[Wiest(2025a)]{wiest2025substrate}
Wiest, M.\,C.
\newblock A quantum microtubule substrate of consciousness is experimentally
  supported and solves the binding and epiphenomenalism problems.
\newblock \emph{Neuroscience of Consciousness}, 2025.
\newblock DOI: \href{https://doi.org/10.1093/nc/niaf011}{10.1093/nc/niaf011}

\bibitem[Wiest(2025b)]{wiest2025spintronic}
Wiest, M.\,C.
\newblock Consciousness and spintronic coherence in microtubules.
\newblock \emph{Communicative \& Integrative Biology}, 18(1), 2025.
\newblock DOI: \href{https://doi.org/10.1080/19420889.2025.2576334}{10.1080/19420889.2025.2576334}

\bibitem[Tononi(2004)]{tononi2004}
Tononi, G.
\newblock An information integration theory of consciousness.
\newblock \emph{BMC Neuroscience}, 5, 42, 2004.

\bibitem[Dehaene and Changeux(2011)]{dehaene2011}
Dehaene, S.; Changeux, J.-P.
\newblock Experimental and theoretical approaches to conscious processing.
\newblock \emph{Neuron}, 70(2), 200--227, 2011.

\bibitem[Friston(2010)]{friston2010}
Friston, K.
\newblock The free-energy principle: a unified brain theory?
\newblock \emph{Nature Reviews Neuroscience}, 11(2), 127--138, 2010.
\newblock DOI: \href{https://doi.org/10.1038/nrn2787}{10.1038/nrn2787}

\bibitem[Rao and Ballard(1999)]{raoballard1999}
Rao, R.\,P.\,N.; Ballard, D.\,H.
\newblock Predictive coding in the visual cortex: a functional interpretation
  of some extra-classical receptive-field effects.
\newblock \emph{Nature Neuroscience}, 2(1), 79--87, 1999.
\newblock DOI: \href{https://doi.org/10.1038/4580}{10.1038/4580}

\bibitem[Jaeger(2001)]{jaeger2001}
Jaeger, H.
\newblock The ``echo state'' approach to analysing and training recurrent
  neural networks.
\newblock GMD Report 148, German National Research Center for Information
  Technology, 2001.

\bibitem[Cleary(2008)]{cleary2008}
Cleary, A.M. (2008).
Recognition Memory, Familiarity, and Déjà vu Experiences.
\textit{Current Directions in Psychological Science}, 17(5), 353--357.

\bibitem[Illman et al.(2012)]{illman2012}
Illman, N.A., Butler, C.R., Souchay, C., \& Moulin, C.J.A. (2012).
Déjà Experiences in Temporal Lobe Epilepsy.
\textit{Epilepsy Research and Treatment}, 2012, 539567.

\bibitem[Martin et al.(2021)]{martin2021}
Martin, C.B., Mirsattari, S.M., Pruessner, J.C., et al.\ (2021).
Relationship between déjà vu experiences and recognition-memory impairments
in temporal-lobe epilepsy.
\textit{Memory}, 29(7), 884--894.

\bibitem[Pe\v{s}lov\'{a} et al.(2018)]{peslova2018}
Pe\v{s}lov\'{a}, E., Mare\v{c}ek, R., Shaw, D.J., et al.\ (2018).
Hippocampal involvement in nonpathological déjà vu:
Subfield vulnerability rather than temporal lobe epilepsy equivalent.
\textit{Brain and Behavior}, 8(7), e00996.

\end{thebibliography}

\end{document}
